\documentclass[11pt]{article}

\usepackage[margin=1in]{geometry}
% Shared preamble for the urcu-txn paper series.
% Included by each paper's main.tex AFTER \documentclass.
%
% Deliberately targets what arXiv's AutoTeX accepts:
%   - pdflatex (arXiv default is TeX Live 2025, matching our local install)
%   - natbib + bibtex, with the generated .bbl committed alongside the source
%     (biblatex .bbl is version-locked to the biblatex/biber pair; natbib's
%     is not, so natbib is the lower-risk path even though arXiv now runs
%     both backends)
%   - pgfplots reading CSV directly, so plot data ships inside the submission
%     and no external plotting toolchain is needed to rebuild a figure

\usepackage[T1]{fontenc}
\usepackage[utf8]{inputenc}
\usepackage{lmodern}
\usepackage{microtype}
\usepackage{xspace}
\usepackage{xcolor}

% This series is dense with long, unbreakable C identifiers
% (cds_list_for_each_entry_reverse, -DURCU_TXN_SW_EXCL_VALIDATE), which TeX
% cannot hyphenate and will happily shove into the margin. emergencystretch
% lets it loosen interword spacing as a last resort instead.
%
% Why not hyphenate typewriter text (\usepackage[htt]{hyphenat}): a hyphen
% inside a C identifier is ambiguous -- the reader cannot tell whether it is
% part of the name. A slightly loose line is a much cheaper defect than an
% identifier the reader has to guess at, in a paper whose point is enablement.
\setlength{\emergencystretch}{3em}

\usepackage{amsmath}
\usepackage{amssymb}
\usepackage{amsthm}

\usepackage{booktabs}
\usepackage{multirow}
\usepackage{graphicx}
\usepackage{subcaption}
\usepackage{siunitx}

\usepackage{tikz}
% positioning: `above=2mm of foo`. Without it TikZ parses "of" as a math
% operator and dies with an opaque PGF Math Error rather than a missing-library
% message, so load it up front.
\usetikzlibrary{positioning,arrows.meta,calc,fit,backgrounds}
\usepackage{pgfplots}
\pgfplotsset{compat=1.18}
\usepgfplotslibrary{groupplots}

\usepackage[ruled,vlined,linesnumbered]{algorithm2e}
\usepackage{listings}
\usepackage{etoolbox}

\usepackage[numbers,sort&compress]{natbib}

% hyperref late, cleveref last.
\usepackage[colorlinks=true,linkcolor=black,citecolor=black,urlcolor=blue]{hyperref}
\usepackage{cleveref}

% ---------------------------------------------------------------------------
% C listings style, for the API and pseudocode excerpts. Enablement is the
% point of a defensive publication, so code excerpts are first-class content
% and must stay legible in print.
% ---------------------------------------------------------------------------
\lstdefinestyle{c}{
  language=C,
  basicstyle=\ttfamily\small,
  keywordstyle=\bfseries,
  commentstyle=\itshape\color{gray!70!black},
  stringstyle=\color{black},
  numbers=none,
  frame=none,
  columns=fullflexible,
  keepspaces=true,
  showstringspaces=false,
  breaklines=true,
  captionpos=b,
  morekeywords={uint64_t,uintptr_t,bool,struct,static,inline,_Atomic},
}
\lstset{style=c}

% ---------------------------------------------------------------------------
% Listings must not split across a page boundary. The excerpts in this series
% are short (a handful of lines each) and are read as single units -- a
% __list_del() whose two stores land on different pages defeats the entire
% point of showing them together. listings has no option for this, so wrap
% every listing in an unbreakable box. If a listing ever grows longer than a
% page this will overfull rather than break: that is the correct failure, since
% such a listing wants to be a float with a caption instead.
% ---------------------------------------------------------------------------
\BeforeBeginEnvironment{lstlisting}{\par\noindent\begin{minipage}{\linewidth}}
\AfterEndEnvironment{lstlisting}{\end{minipage}\par}

\theoremstyle{definition}
\newtheorem{definition}{Definition}
\theoremstyle{plain}
\newtheorem{theorem}{Theorem}
\newtheorem{lemma}{Lemma}
\newtheorem{invariant}{Invariant}
\theoremstyle{remark}
\newtheorem{property}{Property}

% Shared macros for the urcu-txn paper series.
% Included by each paper's main.tex via % Shared macros for the urcu-txn paper series.
% Included by each paper's main.tex via % Shared macros for the urcu-txn paper series.
% Included by each paper's main.tex via \input{../common/macros}.
%
% Keep this class-agnostic: nothing here may assume `article`, so that a paper
% can be retargeted to acmart/IEEEtran by changing only its own preamble.

% ---------------------------------------------------------------------------
% Terminology.
%
% These exist so the series says the same thing the same way every time, and so
% that a terminology decision is changed in one place. The guarantee wording is
% load-bearing: this engine is NOT STM and is NOT serializable. See
% \pseudotxn below and the note in README.md.
% ---------------------------------------------------------------------------
\newcommand{\pseudotxn}{pseudo-transaction\xspace}
\newcommand{\pseudotxns}{pseudo-transactions\xspace}
\newcommand{\mcas}{MCAS\xspace}
\newcommand{\kcas}{$k$-CAS\xspace}
\newcommand{\rcu}{RCU\xspace}
\newcommand{\QSBR}{QSBR\xspace}
\newcommand{\urcu}{userspace RCU\xspace}
\newcommand{\sw}{single-writer\xspace}
\newcommand{\mw}{multi-writer\xspace}
\newcommand{\fliplatch}{flip-latch\xspace}
\newcommand{\ryw}{read-your-own-writes\xspace}

% Engine identifiers.
\newcommand{\urcutxn}{\texttt{urcu-txn}\xspace}
\newcommand{\rcumcas}{\texttt{rcu-mcas}\xspace}
\newcommand{\rcutxnsw}{\texttt{rcu-txn-sw}\xspace}

% ---------------------------------------------------------------------------
% Code and identifiers.
% ---------------------------------------------------------------------------
% \code offers TeX no place to break, and this series is dense with identifiers
% long enough to overflow a line on their own (cds_list_for_each_entry_reverse,
% -DURCU_TXN_SW_EXCL_VALIDATE). Two of them in one sentence is enough to push
% text into the margin no matter how the surrounding prose is stretched.
%
% Since every such identifier already writes its underscores as \_, redefining
% \_ locally gives TeX a legal breakpoint after each one -- silently, with NO
% hyphen. That last part is the whole reason not to reach for hyphenat's [htt]
% option instead: a hyphen inside a C identifier is ambiguous, and the reader
% cannot tell whether it belongs to the name. For a paper whose purpose is
% enablement, an identifier someone has to guess at is a real defect; a slightly
% loose line is not.
%
% Scoped to \texttt's own group, so \_ keeps its normal meaning everywhere else.
%
% \let, not \renewcommand, and \DeclareRobustCommand, not \newcommand: \code is
% used inside \caption, which is a MOVING argument (written to the .aux and
% re-read). A fragile \code expands there and the \renewcommand form dies with
% "Argument of \caption@ydblarg has an extra }" -- a fatal error whose message
% names neither \code nor the caption it came from.
\newcommand{\codeunderscore}{\textunderscore\allowbreak}
\DeclareRobustCommand{\code}[1]{\texttt{\let\_\codeunderscore #1}}
\DeclareRobustCommand{\fn}[1]{\code{#1()}}

% A record in a transaction's write set: {slot, old, new}.
\newcommand{\record}[3]{$\{#1,\,#2 \rightarrow #3\}$}
% An edit on one slot, old to new.
\newcommand{\edit}[2]{$#1 \rightarrow #2$}

% Transaction status values.
\newcommand{\UNDECIDED}{\textsc{undecided}\xspace}
\newcommand{\SUCCEEDED}{\textsc{succeeded}\xspace}
\newcommand{\FAILED}{\textsc{failed}\xspace}
\newcommand{\ABORT}{\textsc{abort}\xspace}

% ---------------------------------------------------------------------------
% Editorial markers. Define \draftmode in the preamble to make these visible;
% they vanish in a submission build. Never ship with \draftmode set.
% ---------------------------------------------------------------------------
\newif\ifdraftmode
\draftmodefalse

% These use \color/\bfseries SWITCHES inside a group rather than
% \textcolor{}{}/\textbf{}, whose arguments are short: an editorial note
% spanning a blank line would otherwise die with "Paragraph ended before
% \@textcolor was complete". Notes routinely span paragraphs, so the switch
% form is the only one that survives real use.
\newcommand{\marker}[3]{\ifdraftmode{\color{#1}\bfseries[#2: #3]}\fi}

\newcommand{\TODO}[1]{\marker{red}{TODO}{#1}}
\newcommand{\NOTE}[1]{\marker{blue}{NOTE}{#1}}
% Claims needing a citation before submission.
\newcommand{\NEEDCITE}[1]{\marker{orange}{CITE}{#1}}
% Numbers currently sourced from source comments rather than a reproducible run.
\newcommand{\UNVERIFIED}[1]{\marker{purple}{UNVERIFIED}{#1}}
.
%
% Keep this class-agnostic: nothing here may assume `article`, so that a paper
% can be retargeted to acmart/IEEEtran by changing only its own preamble.

% ---------------------------------------------------------------------------
% Terminology.
%
% These exist so the series says the same thing the same way every time, and so
% that a terminology decision is changed in one place. The guarantee wording is
% load-bearing: this engine is NOT STM and is NOT serializable. See
% \pseudotxn below and the note in README.md.
% ---------------------------------------------------------------------------
\newcommand{\pseudotxn}{pseudo-transaction\xspace}
\newcommand{\pseudotxns}{pseudo-transactions\xspace}
\newcommand{\mcas}{MCAS\xspace}
\newcommand{\kcas}{$k$-CAS\xspace}
\newcommand{\rcu}{RCU\xspace}
\newcommand{\QSBR}{QSBR\xspace}
\newcommand{\urcu}{userspace RCU\xspace}
\newcommand{\sw}{single-writer\xspace}
\newcommand{\mw}{multi-writer\xspace}
\newcommand{\fliplatch}{flip-latch\xspace}
\newcommand{\ryw}{read-your-own-writes\xspace}

% Engine identifiers.
\newcommand{\urcutxn}{\texttt{urcu-txn}\xspace}
\newcommand{\rcumcas}{\texttt{rcu-mcas}\xspace}
\newcommand{\rcutxnsw}{\texttt{rcu-txn-sw}\xspace}

% ---------------------------------------------------------------------------
% Code and identifiers.
% ---------------------------------------------------------------------------
% \code offers TeX no place to break, and this series is dense with identifiers
% long enough to overflow a line on their own (cds_list_for_each_entry_reverse,
% -DURCU_TXN_SW_EXCL_VALIDATE). Two of them in one sentence is enough to push
% text into the margin no matter how the surrounding prose is stretched.
%
% Since every such identifier already writes its underscores as \_, redefining
% \_ locally gives TeX a legal breakpoint after each one -- silently, with NO
% hyphen. That last part is the whole reason not to reach for hyphenat's [htt]
% option instead: a hyphen inside a C identifier is ambiguous, and the reader
% cannot tell whether it belongs to the name. For a paper whose purpose is
% enablement, an identifier someone has to guess at is a real defect; a slightly
% loose line is not.
%
% Scoped to \texttt's own group, so \_ keeps its normal meaning everywhere else.
%
% \let, not \renewcommand, and \DeclareRobustCommand, not \newcommand: \code is
% used inside \caption, which is a MOVING argument (written to the .aux and
% re-read). A fragile \code expands there and the \renewcommand form dies with
% "Argument of \caption@ydblarg has an extra }" -- a fatal error whose message
% names neither \code nor the caption it came from.
\newcommand{\codeunderscore}{\textunderscore\allowbreak}
\DeclareRobustCommand{\code}[1]{\texttt{\let\_\codeunderscore #1}}
\DeclareRobustCommand{\fn}[1]{\code{#1()}}

% A record in a transaction's write set: {slot, old, new}.
\newcommand{\record}[3]{$\{#1,\,#2 \rightarrow #3\}$}
% An edit on one slot, old to new.
\newcommand{\edit}[2]{$#1 \rightarrow #2$}

% Transaction status values.
\newcommand{\UNDECIDED}{\textsc{undecided}\xspace}
\newcommand{\SUCCEEDED}{\textsc{succeeded}\xspace}
\newcommand{\FAILED}{\textsc{failed}\xspace}
\newcommand{\ABORT}{\textsc{abort}\xspace}

% ---------------------------------------------------------------------------
% Editorial markers. Define \draftmode in the preamble to make these visible;
% they vanish in a submission build. Never ship with \draftmode set.
% ---------------------------------------------------------------------------
\newif\ifdraftmode
\draftmodefalse

% These use \color/\bfseries SWITCHES inside a group rather than
% \textcolor{}{}/\textbf{}, whose arguments are short: an editorial note
% spanning a blank line would otherwise die with "Paragraph ended before
% \@textcolor was complete". Notes routinely span paragraphs, so the switch
% form is the only one that survives real use.
\newcommand{\marker}[3]{\ifdraftmode{\color{#1}\bfseries[#2: #3]}\fi}

\newcommand{\TODO}[1]{\marker{red}{TODO}{#1}}
\newcommand{\NOTE}[1]{\marker{blue}{NOTE}{#1}}
% Claims needing a citation before submission.
\newcommand{\NEEDCITE}[1]{\marker{orange}{CITE}{#1}}
% Numbers currently sourced from source comments rather than a reproducible run.
\newcommand{\UNVERIFIED}[1]{\marker{purple}{UNVERIFIED}{#1}}
.
%
% Keep this class-agnostic: nothing here may assume `article`, so that a paper
% can be retargeted to acmart/IEEEtran by changing only its own preamble.

% ---------------------------------------------------------------------------
% Terminology.
%
% These exist so the series says the same thing the same way every time, and so
% that a terminology decision is changed in one place. The guarantee wording is
% load-bearing: this engine is NOT STM and is NOT serializable. See
% \pseudotxn below and the note in README.md.
% ---------------------------------------------------------------------------
\newcommand{\pseudotxn}{pseudo-transaction\xspace}
\newcommand{\pseudotxns}{pseudo-transactions\xspace}
\newcommand{\mcas}{MCAS\xspace}
\newcommand{\kcas}{$k$-CAS\xspace}
\newcommand{\rcu}{RCU\xspace}
\newcommand{\QSBR}{QSBR\xspace}
\newcommand{\urcu}{userspace RCU\xspace}
\newcommand{\sw}{single-writer\xspace}
\newcommand{\mw}{multi-writer\xspace}
\newcommand{\fliplatch}{flip-latch\xspace}
\newcommand{\ryw}{read-your-own-writes\xspace}

% Engine identifiers.
\newcommand{\urcutxn}{\texttt{urcu-txn}\xspace}
\newcommand{\rcumcas}{\texttt{rcu-mcas}\xspace}
\newcommand{\rcutxnsw}{\texttt{rcu-txn-sw}\xspace}

% ---------------------------------------------------------------------------
% Code and identifiers.
% ---------------------------------------------------------------------------
% \code offers TeX no place to break, and this series is dense with identifiers
% long enough to overflow a line on their own (cds_list_for_each_entry_reverse,
% -DURCU_TXN_SW_EXCL_VALIDATE). Two of them in one sentence is enough to push
% text into the margin no matter how the surrounding prose is stretched.
%
% Since every such identifier already writes its underscores as \_, redefining
% \_ locally gives TeX a legal breakpoint after each one -- silently, with NO
% hyphen. That last part is the whole reason not to reach for hyphenat's [htt]
% option instead: a hyphen inside a C identifier is ambiguous, and the reader
% cannot tell whether it belongs to the name. For a paper whose purpose is
% enablement, an identifier someone has to guess at is a real defect; a slightly
% loose line is not.
%
% Scoped to \texttt's own group, so \_ keeps its normal meaning everywhere else.
%
% \let, not \renewcommand, and \DeclareRobustCommand, not \newcommand: \code is
% used inside \caption, which is a MOVING argument (written to the .aux and
% re-read). A fragile \code expands there and the \renewcommand form dies with
% "Argument of \caption@ydblarg has an extra }" -- a fatal error whose message
% names neither \code nor the caption it came from.
\newcommand{\codeunderscore}{\textunderscore\allowbreak}
\DeclareRobustCommand{\code}[1]{\texttt{\let\_\codeunderscore #1}}
\DeclareRobustCommand{\fn}[1]{\code{#1()}}

% A record in a transaction's write set: {slot, old, new}.
\newcommand{\record}[3]{$\{#1,\,#2 \rightarrow #3\}$}
% An edit on one slot, old to new.
\newcommand{\edit}[2]{$#1 \rightarrow #2$}

% Transaction status values.
\newcommand{\UNDECIDED}{\textsc{undecided}\xspace}
\newcommand{\SUCCEEDED}{\textsc{succeeded}\xspace}
\newcommand{\FAILED}{\textsc{failed}\xspace}
\newcommand{\ABORT}{\textsc{abort}\xspace}

% ---------------------------------------------------------------------------
% Editorial markers. Define \draftmode in the preamble to make these visible;
% they vanish in a submission build. Never ship with \draftmode set.
% ---------------------------------------------------------------------------
\newif\ifdraftmode
\draftmodefalse

% These use \color/\bfseries SWITCHES inside a group rather than
% \textcolor{}{}/\textbf{}, whose arguments are short: an editorial note
% spanning a blank line would otherwise die with "Paragraph ended before
% \@textcolor was complete". Notes routinely span paragraphs, so the switch
% form is the only one that survives real use.
\newcommand{\marker}[3]{\ifdraftmode{\color{#1}\bfseries[#2: #3]}\fi}

\newcommand{\TODO}[1]{\marker{red}{TODO}{#1}}
\newcommand{\NOTE}[1]{\marker{blue}{NOTE}{#1}}
% Claims needing a citation before submission.
\newcommand{\NEEDCITE}[1]{\marker{orange}{CITE}{#1}}
% Numbers currently sourced from source comments rather than a reproducible run.
\newcommand{\UNVERIFIED}[1]{\marker{purple}{UNVERIFIED}{#1}}


% Draft build: `make draft` drops draft.flag to surface TODO/CITE/NOTE markers.
% Submission builds leave it absent, so the markers vanish.
\IfFileExists{draft.flag}{\draftmodetrue}{}

\title{RCU Pseudo-Transactions and the Single-Writer Flip-Latch:\\
       Atomic Multi-Slot Publication for Read-Copy-Update Data Structures}

\author{Mathieu Desnoyers\\ EfficiOS Inc.\\ \texttt{mathieu.desnoyers@efficios.com}}

\date{\today}

\begin{document}
\maketitle

\begin{abstract}
% ---------------------------------------------------------------------------
% FRAMING RULE FOR THE WHOLE SERIES -- do not soften this in editing.
% The guarantee is: atomic multi-slot write, opt-in per-slot read validation.
% NEVER "serializable", NEVER "opacity", NEVER "STM". "Linearizable" stays
% scoped to the commit itself, never to a read+write transaction. Overclaiming
% here is the single most damaging error available to this paper.
% ---------------------------------------------------------------------------
Read-copy-update (RCU) makes a \emph{single} pointer store atomic to concurrent
readers, and readers pay nothing for it in the absence of updates.
This works only when an update is expressible as one store. Updates that must
change several words coherently---a doubly-linked list's back edge, a skip
list's tower---fall
back to copying a larger region, a grace-period drain, or omitting the
operation. Linux's list does the last: it ships no RCU-safe reverse iterator,
because \code{\_\_list\_del()} publishes its forward edge with a marked store
and its back edge with a plain assignment.

We describe an RCU \pseudotxn: a frozen set of \record{slot}{old}{new} records
made visible by one commit step, together with a \sw realization---the
\fliplatch---in which every proxy in a group names one shared selector word, so
a single store switches the whole set. Two transacted structures follow: a
bidirectional list whose two chains stay mutually coherent, so a reader may walk
it backward, and a hash list.

The guarantee is deliberately weaker than software transactional memory. The
write set is atomic: every record resolves to its old value before the commit
and to its new value after. But a reader is \emph{not} a transaction of any
kind---each of its loads resolves independently, so a traversal spanning several slots may
straddle a commit. There is no opacity and no multi-read snapshot; a reader gets
monotonicity, never new-then-old.

The weakening is by design. Where a structure is mutated in place, isolation is
a read-side cost: state on every element, work on every dereference---and the
read side is what RCU exists to keep free. A \pseudotxn's marker instead rides
in spare bits of the slot itself, the word the reader loads anyway, so it adds
no per-element state and, absent a concurrent commit on that slot, no extra
load. The bulk of the cost moves to the writer, in the width of a commit.
\end{abstract}

% ===========================================================================
\section{Introduction}
\label{sec:intro}

Read-copy-update rests on a contract. Readers pay nothing --- no lock, no
barrier, no write to shared state --- and in exchange a writer gets exactly one
atomic operation: a single pointer store, bracketed by building the new structure
where readers cannot reach it and reclaiming the old one after a grace period.
The contract has held for decades because the structures \rcu serves are
read-dominated, so making the majority free is worth constraining the minority.

The constraint is real, though, and it is one store. An update that must change
several words coherently has no atomic operation to express itself with. The remedies
are familiar and each gives something up: copy the enclosing region so the
change becomes one pointer swap, and pay to copy; drain the readers with a grace
period so nobody spans the inconsistent window, and pay to wait; impose memory
ordering so the reader is steered around the inconsistent window instead of being
made immune to it, and pay in read-side performance~\citep{JoshTriplettPhD}; or
decline the operation.

\Cref{sec:onestore} looks at a library that took the third option, and it is
the cleanest evidence available that the problem is real, and it is not confined
to a niche library. Linux's \rcu list --- the most widely deployed there is ---
publishes only its forward chain: \code{\_\_list\_del()} commits the forward
edge with \code{WRITE\_ONCE()} and writes the back edge with a plain
assignment, and \code{\_\_list\_add\_rcu()} writes it \emph{after} the release
that publishes the node. Userspace \rcu's list has the identical shape. The
consequence is not a subtle ordering nit but a semantic one: the two chains are
changed by two stores and nothing makes them one event, so there is an interval
in which they describe different lists --- and a removed node survives as a
ghost until its grace period elapses, so the \code{prev} view can trail the
\code{next} view by that long. No barrier fixes either problem. So both
libraries ship forward iterators and no \rcu-safe reverse walk, and Linux goes
further, poisoning the back pointer in \code{list\_del\_rcu()} so no concurrent
\rcu reader can follow it. Reverse traversal is not exotic: the same libraries' non-\rcu lists offer it as
a matter of course, so its absence from the \rcu variants marks a gap \rcu could
not close, not a feature nobody wanted.

\subsection*{Contributions}

This paper describes a facility that gives \rcu writers a second atomic operation ---
a set of word-sized slots that change together --- without imposing
significant additional cost on \rcu readers in the absence of updates.
Specifically:

\begin{itemize}
\item \textbf{A \pseudotxn model} (\cref{sec:model}): a frozen set of
      \record{slot}{old}{new} records made visible by one commit step. We state
      what it guarantees, and state what it does not before any evaluation ---
      including the property that a reader is not a transaction of any kind, and
      the weaker
      one that it does have: monotone resolution, never new-then-old
      (\cref{sec:monotone}).

\item \textbf{The \fliplatch} (\cref{sec:fliplatch}): a \sw realization in which
      every proxy in a group names one shared selector word, so a single release
      store switches the set. It has no compare-and-swap, no conflict detection
      and no abort --- and correspondingly requires writer exclusion, which we
      state as a precondition of the mechanism rather than a caveat on it.

\item \textbf{A per-record tag scheme} (\cref{sec:tags}): the marker lives in
      spare bits of the slot rather than in per-element state, which is what
      takes the reader's added loads to zero --- the test is on a register the
      reader already holds, and costs it only a predicted branch. Because a
      \pseudotxn may span structures whose pointers are tagged differently, the
      tag must travel with the record rather than be fixed per translation unit.

\item \textbf{Two transacted structures} (\cref{sec:wrappers}): a bidirectional
      list whose two chains stay mutually coherent --- so the reverse walk
      \code{rculist} declines to offer becomes one a reader may take, within
      the guarantee and the scope \cref{sec:coherence} states --- and a hash list.
      Several edits compose into one commit, including edits that touch the same
      slots, through a read-your-own-writes discipline (\cref{sec:swcomposition});
      we give the two obligations that remain.

\item \textbf{Measurements on the reader axis} (\cref{sec:measured}), scaling to
      192 cores: on the same traversal and a byte-identical element, readers of
      the transacted list are indistinguishable from plain \rcu{} --- within 3\%
      throughout, against a machine drift of 0.8\% --- while the facility also
      delivers the reverse walk plain \rcu cannot offer under mutation at all.
      Against the schemes that \emph{do} offer it, and which give readers
      strictly more, the reader pays 14--21\% for existence and 27--50\% for
      RLU. A single writer with readers running pays about twice, and a profile
      says where that goes --- 3.4\% of cycles in the commit, the rest in
      allocation and grace-period reclaim.

\item \textbf{The design space we did not explore} (\cref{sec:variations}):
      selector encodings, tag placement and width, proxy allocation, settle
      policies, and reclamation schemes, with the cost of each alternative
      rather than only the name of it.
\end{itemize}

\paragraph{Scope.} This paper covers the \sw case, where writers are mutually
excluded by the embedder. A companion paper covers the concurrent facility, which
attacks the same model with a software multi-word compare-and-swap and pays for
contention detection with machinery this one does not need. The two share the
model of \cref{sec:model} and the status vocabulary, deliberately: an embedder
written against one migrates to the other without rewriting its commit handling.
Properties that hold only under exclusion are marked where they appear.

\subsection{Where the cost goes}

Multi-pointer atomicity under \rcu is not a new goal, and the established
answers pay for it in the same place: state attached to every element, and
work on the reader's path. This facility pays in neither. \Cref{sec:costcomparison}
places it against those answers---existence structures, RLU, and STM---in detail,
and \cref{sec:measured} measures both halves of the claim on a list: what the
reader does not pay, and what the writer does.

% ===========================================================================
\section{Background}
\label{sec:background}

This section fixes only the vocabulary the rest of the paper needs. Readers
familiar with \rcu will find nothing new in \cref{sec:rcu} and
\cref{sec:buildpublish}; \cref{sec:onestore} is where the problem is stated, and
it is stated on shipped code rather than in the abstract.

\subsection{Read-copy-update}
\label{sec:rcu}

\rcu~\citep{mckenney1998rcu,desnoyers2012urcu} divides a mutation between two
parties and gives each of them a different contract. A reader delimits its
traversal with \code{rcu\_read\_lock()} and \code{rcu\_read\_unlock()}, marking an
\rcu \emph{read-side critical section}. On the
common implementations these are not locks at all: they exclude nothing and
contend with nothing, and in \urcu's \QSBR mode
\code{rcu\_read\_lock()} is an empty function --- a debug assertion and nothing
else. What a reader gets in exchange is narrow but sufficient: any object it
reaches inside the critical section remains allocated until that section ends.

That is not magic, and it is worth saying where the cost went rather than
leaving an empty function to look like a free lunch. \QSBR buys its empty read
lock by moving the reader's obligation out of the critical section and onto the
thread: a
registered thread must periodically call \code{rcu\_quiescent\_state()} to
announce that it holds no \rcu-protected references, and must call
\code{rcu\_thread\_offline()} before blocking for any length of time. The
announcement is cheap --- it compares the global grace-period counter against
the thread's own and returns immediately when they agree --- but somebody has to
make it, and the writer's grace period cannot end until every registered thread
has. So the cost is not eliminated; it is amortized across a whole thread's
worth of read-side critical sections and relocated to a point the application
chooses.
The price is structural rather than instruction-level: the application must
\emph{have} such points and must be written to reach them promptly, which is why
\QSBR is at once the fastest flavour and the most intrusive to adopt. The other
flavours make the opposite trade, paying a little inside the critical section so
the application owes nothing between them.

That promise is kept by the writer, not by the reader, and it is kept by
waiting. A \emph{grace period} is an interval that every pre-existing read-side
critical section has been observed to leave. Once one elapses, no reader can still hold a
reference obtained before it began, so anything unlinked before it began is
unreachable and may be freed. The writer waits synchronously
(\code{synchronize\_rcu()}) or defers a callback (\code{call\_rcu()}); either
way the cost of reclamation is charged entirely to the side doing the mutating.

Publication is the mirror image. A writer makes an object reachable with
\code{rcu\_assign\_pointer()}, a release store, and a reader reaches it with
\code{rcu\_dereference()}, which returns the pointer and guarantees that the
loads dependent on it see the object's initialized contents rather than
pre-initialization garbage. On every architecture Userspace \rcu supports but one, the
reader's half of that is free: the dependency between loading a pointer and
loading through it is respected by the hardware, so \code{rcu\_dereference()}
emits an ordinary load and no barrier.\footnote{The one exception is Alpha, whose
hardware does not order dependent loads; Userspace \rcu supports it, and enforces
the dependency there with an explicit \code{cmm\_smp\_read\_barrier\_depends()}. Elsewhere the dependency
is respected directly: the C11 \code{consume} ordering meant to express it is
unusable in practice, since compilers promote it to
\code{acquire}~\citep{boehm2008memorymodel,batty2011cppconcurrency}, so Userspace
\rcu relies on the dependency itself.}

The whole arrangement rests on one sentence: \emph{readers pay nothing, and in
exchange a writer gets exactly one atomic operation.} The rest of this paper is about
the second half of that sentence.

\subsection{Build, publish, reclaim}
\label{sec:buildpublish}

The one atomic operation is a single pointer store, and it sits in the middle of a
three-phase discipline that every \rcu mutation follows:

\begin{enumerate}
\item \textbf{Build.} Allocate and fully initialize the new nodes \emph{where no
reader can reach them}. Because they are unreachable, this phase is
single-threaded by construction: no ordering is needed, no barrier is paid, and
an operation that fails here has published nothing and can free its work
immediately with no grace period.
\item \textbf{Publish.} One release store makes the new structure reachable.
This is the linearization point, and it is the only instant at which concurrent
readers change what they see. It is also the paper's pivot: \emph{everything
before it is private, everything after it is committed, and there is exactly one
of it.}
\item \textbf{Reclaim.} The nodes the publish displaced were reachable, so they
may still be held by readers that started earlier. Free them after a grace
period, never before.
\end{enumerate}

The middle phase is the constraint. Build is unconstrained and reclaim is merely
deferred, but publish must be \emph{one store}, because a store is the only
thing commodity hardware will make atomic with respect to a reader that has taken
no lock and paid no barrier.\footnote{On hardware with transactional memory, a
whole multi-word update can be published atomically to a lockless
reader~\citep{ramadan2007txlinux}. Such hardware is not universal; we assume only
the single-word store and compare-and-swap that hardware provides everywhere.}

This frame carries through the whole series, and it is worth stating the
punchline now so that \cref{sec:model} lands as an extension rather than a
replacement: a \pseudotxn does not add a fourth phase. It \emph{widens the
second one}. Build and reclaim keep their meanings exactly; commit is the
publish step, still a single store, and what changes is how many slots that one
store moves.

\subsection{Why one store is not enough}
\label{sec:onestore}

An update that must re-point two edges coherently has no atomic operation to express
itself with, and the consequence is not hypothetical. It is visible in the
doubly-linked list of the most widely deployed \rcu implementation there is: the
one in the Linux kernel~\citep{linux2026rculist}.

Removing a node \code{E} from between \code{P} and \code{N} re-points two
reader-visible edges, and the kernel writes each with a different kind of store:

\begin{lstlisting}
static inline void __list_del(struct list_head *prev, struct list_head *next)
{
        next->prev = prev;              /* N->prev: E -> P   plain store  */
        WRITE_ONCE(prev->next, next);   /* P->next: E -> N   marked store */
}
\end{lstlisting}

The forward edge gets \code{WRITE\_ONCE()}. The back edge gets a bare \code{=},
to a node readers can reach at that very moment. Userspace \rcu's
\code{cds\_list\_del\_rcu()} has the identical shape, down to the ordering of
the two lines. And the asymmetry is not an accident of the delete path ---
\code{\_\_list\_add\_rcu()} makes it starker still:

\begin{lstlisting}
new->next = next;                              /* private: unreachable */
new->prev = prev;                              /* private: unreachable */
rcu_assign_pointer(list_next_rcu(prev), new);  /* release: publishes new */
next->prev = new;                              /* plain store, AFTER publish */
\end{lstlisting}

The forward edge is published with a release, and the back edge is written
\emph{afterwards}, with a plain store, to a node that is already reachable.
Between those two lines there is an interval in which \code{prev->next == new}
while \code{next->prev == prev}: the forward chain contains the new node and the
backward chain skips it. This needs no weak-memory argument --- it is plain
program order, so the window is real on \emph{every} architecture.

The kernel is candid about the consequence. \code{list\_del\_rcu()} does not
merely leave the back pointer unmaintained; it destroys it,
\code{entry->prev = LIST\_POISON2}, and the \rcu list ships forward iterators
only. Userspace \rcu draws the same line by omission: its plain list offers
\code{cds\_list\_for\_each\_prev()} and
\code{cds\_list\_for\_each\_entry\_reverse()}, and its \rcu list offers neither.
The backward chain exists in both libraries and every writer maintains it. A
reader may not walk it.

\paragraph{The kernel's narrow exception proves the rule.} In December 2024 the
kernel gained a way to read a \code{prev} pointer under
\rcu~\citep{brauner2024listbidir}: \code{list\_bidir\_del\_rcu()}, which skips
the poisoning, plus \code{list\_bidir\_prev\_rcu()}, an accessor. It was
suggested and reviewed by McKenney himself, so it is a considered answer rather
than an oversight being patched --- and it is worth being precise about what it
does and does not deliver, because it is the closest thing to a counterexample
this paper has.

It does not make the two chains agree. \code{\_\_list\_del()} still writes the
back edge with a plain store, and nothing orders that store against the forward
edge's \code{WRITE\_ONCE()}; \code{\_\_list\_add\_rcu()} still writes it after
the release. What the patch buys is that a deleted node keeps its own \code{prev}
instead of having it poisoned, so a reader standing on a ghost can step backward
out of it. That is a genuinely useful thing and it is not a coherent reverse
walk. Consistent with that, the kernel added no reverse iterator, and the API
costs a whole-list mode: \code{list\_del\_rcu()} and \code{list\_bidir\_del\_rcu()}
may not be used on the same list. Its sole in-tree caller as of
\code{next-20260226}, \code{\_\_ns\_tree\_adjoined\_rcu()} in
\code{kernel/nstree.c}, takes exactly one adjacency hop --- \code{prev}
\emph{or} \code{next} --- and never compares the two directions against each
other. It never needs coherence, so it never misses it.

\paragraph{Ordering is not what is missing.} It is tempting to read all of this
as an ordering bug: publish \code{prev} with \code{rcu\_assign\_pointer()}, read
it with \code{rcu\_dereference()}, and the reverse walk falls out. It does not,
and seeing why is the point of this subsection. Suppose both edges were properly
published. They are still \emph{two} release stores, hence two publication
instants, and a reader that arrives between them sees \code{P->next == N} while
\code{N->prev == E}. The forward chain has dropped \code{E}; the backward chain
still contains it. Neither view is corrupt and neither is torn. They simply
describe different lists, and a reader that walks forward and then hops backward
--- the operation a reverse iterator exists to serve --- gets an answer no
sequential list could ever produce.

What is missing is \emph{atomicity across two slots}, and \rcu has no way to ask
for it. That is why the fix has to be a mechanism rather than a barrier, and it
is why two libraries, independently, chose to withhold the operation instead.

There is a second, subtler failure that survives even careful ordering, and it
is specific to \rcu. An unlinked node is a ghost: it must remain allocated for a
grace period, and it keeps its own \code{next} and \code{prev} pointing at its
former neighbours. The live forward chain drops it at the instant of the
publish. A reverse walker stepping through \code{prev} pointers can keep landing
on it for as long as the ghost survives. The backward view therefore trails the
forward view by up to a grace period --- not by a few instructions, but by the
full reclamation latency of the algorithm. No amount of barrier placement closes
a window whose width is a grace period.

So the doubly-linked list, the textbook structure, is already past what one
store can express, and the library's response was to remove the operation from
the interface. \Cref{sec:fliplatch} updates both edges atomically, which
\cref{sec:wrappers} uses to give the reverse walk back.

% ===========================================================================
\section{The Pseudo-Transaction Model}
\label{sec:model}

% ---------------------------------------------------------------------------
% SCOPE NOTE (drafting): this section is deliberately FACILITY-AGNOSTIC. It must
% define only what the single-writer flip-latch (Sec. 4) and the multi-writer
% MCAS facility (companion paper) BOTH realize, because the companion paper
% inherits these definitions verbatim. Anything that is a property of one
% realization -- descriptors, install/settle, the tri-state status word, abort
% -- belongs there, not here. In particular there is no abort in this model:
% under writer exclusion a commit cannot fail for contention, and the ABORT
% status exists only because the two front-ends share a status enum.
% ---------------------------------------------------------------------------

\subsection{Records and pseudo-transactions}

\begin{definition}[Record]
\label{def:record}
A \emph{record} is a triple \record{s}{o}{n}, where $s$ is the address of a
word-sized \emph{slot}, $o$ is the value the slot must hold at the commit
point,\footnote{Why \emph{at}, not \emph{up to} the commit point? See
\cref{sec:nonguarantees}.} and $n$ is the value it takes afterwards.
\end{definition}

\begin{definition}[Pseudo-transaction]
\label{def:pseudotxn}
A \pseudotxn is a set of records over pairwise-distinct slots that is
\emph{frozen} --- closed to further records --- before a single \emph{commit
step} makes the whole set take effect at once.
\end{definition}

Two properties of \cref{def:pseudotxn} carry the entire model, and they concern
different things: \emph{what} takes effect, and \emph{when}. They also happen at
different times, and the order matters.

\emph{Freezing} settles the \emph{what}, and it comes first. The set is closed
before the commit step, not at it: once the commit procedure is under way the
membership is fixed, and no record may still join. So what will become visible
is decided in advance rather than accreting as the commit proceeds.

The \emph{single step} settles the \emph{when}, and it comes last. The commit
step is one store, not a protocol, so the already-frozen set takes effect all at
once --- there is no interval in which part of it is visible and part is not.

These are two distinct moments in sequence, and only the second is the commit
step. Staging the set is what closes it; the single store that follows is what
makes the closed set visible; and any work that tidies up afterwards changes
nothing a reader can see (\cref{prop:lp}). The staging and the tidying are not
part of the commit step and are not linearization points --- a fact easier to
see once the mechanism is concrete (\cref{sec:lifecycle}), where freezing
coincides with the start of installation and the commit step is the lone store
that flips the set from old to new.

\subsection{A transaction in the filesystem sense}
\label{sec:fssense}

The word \emph{transaction} carries two meanings, and only one of them applies
here. A journaling filesystem stages a batch of block writes, closes the batch,
and writes a commit record; before that record the old state is authoritative,
after it the new state is, and no partial batch is ever authoritative. It calls
this a transaction, and it has never meant isolation by it. A database or an
STM means something else: a unit that is atomic \emph{and} isolated, where
isolation constrains what concurrent \emph{readers} may observe. A journaling
filesystem can offer that second guarantee too, under a different name --- a
snapshot, a point-in-time view held consistent against ongoing writes --- and it
is well known to be expensive, in copy-on-write traffic and in space. Isolation
is not cheap anywhere; the only question is whether the workload needs it.

Our model is the first of these, essentially point for point:

\begin{table}[tb]
\centering
\small
\begin{tabular}{@{}ll@{}}
\toprule
\textbf{Journaling filesystem} & \textbf{Pseudo-transaction} \\
\midrule
journal transaction, closed      & record set, frozen \\
journal blocks                   & \record{s}{o}{n} records \\
commit record                    & the commit step \\
before commit: old state         & before commit: every record resolves $o$ \\
after commit: new state          & after commit: every record resolves $n$ \\
never a partial batch            & no slot, and no instant, is ever torn \\
readers are not transactional    & readers are plain \rcu readers \\
gives atomicity, not isolation   & gives atomicity, not opacity \\
\bottomrule
\end{tabular}
\caption{The model is the journaling sense of \emph{transaction}, not the
database sense. The staged record set is the journal; the commit step is the
commit record.}
\label{tab:fssense}
\end{table}

This is why we say \pseudotxn rather than transaction. The prefix is not an
apology for an incomplete STM; it is a disambiguation for readers who take
\emph{transaction} to imply isolation. What the prefix marks as absent is
isolation, and the reason it is absent is the subject of
\cref{sec:whyweakening}: isolation is paid for on the read side, and the \rcu
read side is precisely what this facility keeps free.
% -----------------------------------------------------------------------
% NOTE (drafting; retained deliberately -- see README framing rule).
% If space is tight this table is the first thing to cut to prose --- but
% the point it makes is not optional. It is what stops a concurrency reader
% filing the model under "weak STM" in the first thirty seconds.
% -----------------------------------------------------------------------

\subsection{What is guaranteed}
\label{sec:guarantees}

\begin{property}[Write-set atomicity]
\label{prop:atomicity}
Every record in a committed \pseudotxn resolves to its old value $o$ before the
commit step and to its new value $n$ after it. One step switches the whole
frozen set.
\end{property}

\begin{property}[No tearing]
\label{prop:notearing}
No slot, and no instant, is ever torn: at every instant, every slot holds
either a value that resolves to $o$ or one that resolves to $n$, and never
anything else.
\end{property}

\begin{property}[The commit linearizes at one point]
\label{prop:lp}
The commit is a linearizable \emph{write operation}: it takes effect atomically
at a single instant, the commit step, which is its lone linearization
point~\citep{herlihy1990linearizability}. The staging that precedes it and the
tidying that follows change no observable value, so neither is a linearization
point.
\end{property}

The scope of \cref{prop:lp} is the write, and only the write --- and the
distinction is the reason for the \emph{pseudo-} prefix, not pedantry. A single-slot
read is itself a linearizable operation, taking effect at its own load. What is
\emph{not} a linearizable operation is a reader's traversal of several slots: it
is a sequence of loads, not one operation, so it has no linearization point of
its own and may straddle the commit (\cref{sec:nonguarantees}). Saying the write
linearizes is therefore a precise and limited claim. It is not the claim that a
combined read-and-write transaction is linearizable, and still less that it is
serializable --- neither holds, and neither is asserted anywhere in this paper.

% -----------------------------------------------------------------------
% NOTE (drafting; retained deliberately -- see README framing rule).
% Keep "linearizable" scoped to the WRITE, exactly as in Property 3. The
% facility's own words: "This linearizes the WRITE." Never extend it to a
% read-plus-write bracket -- that bracket is not atomic and claiming
% otherwise is the single most damaging error available to this series.
% -----------------------------------------------------------------------

\begin{property}[Cross-structure atomicity]
\label{prop:crossstructure}
Because the record set may span data structures, the commit is one
linearization point \emph{across all of them}. This is strictly stronger than
composing independent \rcu structures, where a node can be reachable in $A$
before it is reachable in $B$. The unit of cross-structure atomicity is exactly
one \pseudotxn: two commits are two linearization points, the same as two
unrelated \rcu structures.
\end{property}

\paragraph{Worked example.} Consider a cache with two access paths: entries are
found by key through a hash index, and ordered by recency through a linked list.
Both are \rcu structures and both are read without a lock. Inserting an entry
$E$ has to make it findable and make it ordered.

Plain \rcu forces the writer to choose an order, and neither choice is right.
Publish into the index first, and there is an interval in which a reader finds
$E$ by key, follows it, and discovers it is in no position in the recency list.
Splice into the list first, and there is an interval in which $E$ occupies a
position in the recency list while a lookup for its key still reports it absent.
The interval is not a few instructions --- the writer can be preempted between
the two stores --- and no amount of barrier placement removes it, because there
genuinely are two publications. Callers are left to tolerate the skew, which
usually means the structures grow a validity flag, or a lock, and the reason
\rcu was chosen erodes.

A \pseudotxn names the slots of both structures in one record set: the index
bucket slot that must come to name $E$, and the list edges that splice it into
position. One commit publishes all of them. $E$ becomes findable and ordered at
the same instant, and the question of which to publish first stops existing,
because there is only one publish.

\Cref{prop:crossstructure} also fixes the boundary, and it is a real one. The
unit of cross-structure atomicity is one \pseudotxn, not one writer and not one
critical section. An operation that commits the index edit and then commits the
list edit has produced two linearization points and is back to exactly the skew
above --- the facility has made the composition \emph{available}, not automatic.
Folding edits into one commit is sound even when they touch the same slots: the
facility reconciles a writer's edits against its own pending ones through
read-your-own-writes (\cref{sec:swcomposition}). Cross-structure composition,
where the slots cannot coincide, is the easy case of that.

\subsection{What is not guaranteed}
\label{sec:nonguarantees}

We state the limits here, before any evaluation, because they are constitutive
of the model rather than defects in it.

% ---------------------------------------------------------------------------
% GUARDRAIL (do not delete in an editing pass).
%
% "a transaction of ANY KIND" -- never "a reader is not a \pseudotxn". This
% paper said the latter in five places and it is a genuine defect, not a style
% nit: "a reader is not a pseudo-transaction" parses just as naturally as "a
% reader is not MERELY a pseudo-transaction (it gets a real one)", which is the
% exact inverse of the claim and the most damaging misreading the paper admits.
%
% The drift is mechanical and WILL recur: the \pseudotxn macro is an inviting
% substitution for the bare word "transaction", and here the word must stay
% general precisely because the denial has to cover both classes at once. The
% \pseudotxn is a WRITER-side object; denying that a reader is one is nearly
% vacuous. The real claim is that the reader gets no transactional treatment
% whatsoever -- no bracket, no read set, no atomicity across loads.
%
% README's framing rule states it correctly; match it verbatim.
% ---------------------------------------------------------------------------
\textbf{A reader is not a transaction of any kind.} It is a sequence of
single-slot reads, each linearizing at its own access. It does have an \rcu
read-side critical section --- but a critical section is not a transaction. It
keeps reachable objects alive; it does not make the loads inside it a unit.
There is no read set and nothing that begins or commits on the read side: the
\pseudotxn is the writer's object, and a reader meets it one slot at a time. A
traversal spanning several slots may therefore \emph{straddle} a commit: it can
observe the \pseudotxn in a slot it reads after the commit and not in one it
read before.
\Cref{prop:notearing} still holds throughout --- no slot and no instant is torn
--- but a reader's several reads are simply not a mutual snapshot.

\textbf{There is no opacity and no multi-read snapshot.} This is weaker than
STM, deliberately and in exactly one direction: readers are plain \rcu readers
that pay no per-read barrier. An embedder that needs a multi-read snapshot
layers its own versioning on top, as with any value-based \mcas.

\textbf{A record's old is a value at an instant, not a value over an interval.}
\Cref{def:record}'s $o$ is what the slot must hold \emph{at} the commit point,
not what it must have held throughout. An operation whose correctness requires
\emph{detecting} an away-and-back change must carry its own version in the word.
% -----------------------------------------------------------------------
% NOTE (drafting; retained deliberately -- see README framing rule).
% Do NOT write "a validated slot is CHECKED to hold its expected value" in
% P1. The single-writer facility performs no check: it has no install CAS, so
% nothing verifies o against memory. Under exclusion the writer is entitled
% to know, and o is used for RESOLUTION (what a proxy returns pre-flip), not
% for validation. The concurrent facility's install CAS is what turns this
% into a check -- so the checking language belongs in P2, and the model here
% must state the property without implying an enforcement that does not
% exist. Sec. 7.4 explains where this DOES bite in P1: through composition.
% -----------------------------------------------------------------------

\subsection{Monotonicity: never new-then-old}
\label{sec:monotone}

The previous section understates what a reader gets, and the gap matters. A
reader's reads are not mutually atomic, but neither are they unrelated.

\begin{property}[Monotone resolution]
\label{prop:monotone}
The word that decides a \pseudotxn's outcome is written once and never written
back. A reader resolving several slots of one \pseudotxn over time therefore
observes a monotone sequence --- $o \ldots o, n \ldots n$ --- and
\textbf{never $n$ then $o$}. A reader may lag a commit; it never regresses
across one.
\end{property}

\Cref{prop:monotone} is strictly more than treating each slot as an independent
register and strictly less than a snapshot, and that middle ground is
exploitable rather than merely tolerable (\cref{sec:publishorder}).

% -----------------------------------------------------------------------
% NOTE (drafting; retained deliberately -- see README framing rule).
% DO NOT call this "causal". "Causal consistency" names a specific model
% that is WEAKER than linearizability -- it lets observers disagree on the
% order of unrelated writes -- so borrowing the word would make the facility
% read as claiming less than it provides. "Never new-then-old" is precise,
% memorable, collides with no existing model's name, and is the facility's own
% phrasing. Equally, do not write "per-slot linearizability": true of each
% read, but it implies reads are unrelated and so undersells Property 5.
% -----------------------------------------------------------------------

\paragraph{Scope.} \Cref{prop:monotone} has a precondition that must travel
with it. Monotonicity of the deciding word is unconditional: it is one word,
stored once. What a reader's \emph{sequence} of resolutions inherits from it
depends on how the reader got from one slot to the next. A
\emph{dependency-chained} traversal --- each hop's address derived from the
value the previous hop loaded --- carries the order for free on every
architecture \urcu supports, as does any build using the default C11
dereference. A reader that hops \emph{without} that chain, by re-reading a
pointer it cached earlier or by walking between two independently-reached
slots, has no such edge: under \code{-DURCU\_DEREFERENCE\_USE\_VOLATILE} on
weakly-ordered hardware its two loads may be reordered, and it can observe the
new value on the later slot and the old on the earlier one. To avoid it, such a reader must use an acquire load or a dependency.

We state the build condition rather than relegating it, because it is the
difference between a guarantee and a hope. An embedder who misses it writes a
race that no test on a strongly-ordered machine will ever reproduce.

\subsection{Publish order, not snapshots}
\label{sec:publishorder}

\Cref{sec:nonguarantees} invites an obvious objection. If a reader gets no
snapshot, and a traversal touching several slots may straddle a commit, how is
any multi-slot traversal ever correct? The answer is not that snapshots are
recovered somewhere later. It is that a snapshot is only one way of providing what
the reader actually needs, and there is a cheaper way to provide it.

What a snapshot provides is \emph{independence from timing}: it makes a reader's
conclusions hold no matter when its individual loads happened. But a reader does
not read in an arbitrary order. It reads in the order its traversal dictates,
and that order is usually fixed by the structure --- a tree is descended, a list
is walked, a hash bucket is entered at its head. If the order in which edits
\emph{become visible} is made to agree with the order in which the reader
\emph{looks}, timing-independence stops being something to buy at all.

\paragraph{The rule.} Suppose a reader resolves slots in a fixed order
$s_1, s_2, \ldots, s_k$, and the writer splits the edit across several
\pseudotxns rather than one. Then:

\begin{quote}
\emph{Commit in the reverse of the reader's read order: publish $s_k$ first and
$s_1$ last.}
\end{quote}

This ordering rule is, to our knowledge, first stated by
Triplett~\citep{JoshTriplettPhD}.

The consequence is worth spelling out, because it is a short argument with a
disproportionate payoff. Write $t_i$ for the instant $s_i$ is published and
$a_i$ for the instant the reader resolves it. The reader's order gives
$a_1 < a_2 < \cdots < a_k$; the rule gives $t_1 > t_2 > \cdots > t_k$. Now
suppose the reader observes $s_i$ as \emph{new}, i.e.\ $a_i > t_i$. For any
later slot $s_j$ with $j > i$ we have $a_j > a_i > t_i > t_j$, so the reader
observes $s_j$ as new too. Once a reader crosses into the new state it stays
there for the remainder of its traversal. It sees a prefix of old and a suffix
of new, and the crossing point is wherever the commits caught it.

\paragraph{Why a prefix of old and a suffix of new is enough.} The two
observations the reader can be left with are not equally troublesome, and the
asymmetry is the whole point.

Old-then-new is a straddle, not a contradiction. It is precisely what an
ordinary \rcu reader has always tolerated: a reader that begins before a publish
sees the old structure, one that begins after sees the new, and one that
overlaps sees the old up to the point the publish caught it and the new
thereafter. Every deployed \rcu container already asks its callers to accept
this, and callers do, because the alternative --- excluding readers for the
duration of a traversal --- is the cost \rcu exists to avoid.

New-then-old is a different animal. A reader that observes the new value of one
slot and then the old value of another has seen two effects of a single state
transition in reverse order, and no sequential execution of the data structure
produces that. There is nothing the reader can do about it locally: it cannot
tell the stale load from a fresh one, so its only recourse is to re-read
everything and compare --- which is to say, to build a snapshot. \emph{This is
the observation that forces snapshots, and it is the one the model already
forbids.} \Cref{prop:monotone} rules it out by construction, because the
deciding word is written once and never written back.

So the design rule does not weaken the reader's obligations and then apologize
for it. It removes the only observation a reader genuinely could not have coped
with, and leaves behind exactly the straddle that \rcu programmers have been
handling correctly for two decades.

\paragraph{A worked example.} Re-parent a node $E$ from $P$ to $Q$ in a tree
that carries parent pointers. Two reader-visible edges move: the forward edge
($Q$'s child link now names $E$) and the back edge ($E$'s parent link now names
$Q$). A reader descends first and walks back up second, so it reads the forward
edge before the back edge. By the rule, publish them in the opposite order:

\begin{enumerate}
\item Commit the back edge: \code{E->parent}: \edit{P}{Q}.
\item Commit the forward edge: the child links of $P$ and $Q$.
\end{enumerate}

A reader that descends and finds $E$ in its new position under $Q$ read the
forward edge after the second commit, hence after the first, so when it walks
back up it finds \code{E->parent} already naming $Q$, and the two agree. The excluded
observation --- descending into the new position and then being told by the
parent pointer that $E$ still hangs off $P$ --- is exactly the
$(\mathit{new}, \mathit{old})$ pair the ordering makes unreachable. Publishing
the forward edge first would permit it.

The rule keeps the two edges of the move coherent; it does not make the move
atomic with respect to a whole traversal. Moving $E$'s subtree under lockless
readers has the consequences \rcu moves always have: a concurrent reader can skip
the subtree, visit it twice, or --- where a lookup's own path is its key --- find
a key still present at a location it has left. When those are tolerable, and how
a structure that cannot tolerate them contains them, is the subject of later
papers in this series. What the model here secures is only the narrower guarantee
the rule provides: the two edges are never seen to disagree, so no reader
observes the move half-completed.

\paragraph{Why this is safe, and where it stops.} The argument above quietly
assumed both orders it used, and both have to be earned. The writer's side is
free here: under exclusion (\cref{sec:exclusion}) the commits are two release
stores issued in program order by one thread. The reader's side is the one that
needs care, and it is exactly the precondition of \cref{sec:monotone} --- the
reader's loads must be genuinely ordered, not merely written in that sequence.

They are, and for a reason worth noticing rather than assuming: the reader
reaches $E$ by loading the forward edge, and then computes the address of
$E \rightarrow parent$ \emph{from the value that load returned}. The second load
depends on the first. The fixed-read-order idiom is therefore not something that
must be reconciled with the dependency-chaining requirement of
\cref{sec:monotone} --- it \emph{is} a dependency-chained descent, which is why
it is safe as stated and why it costs the reader no barrier.

It is worth being concrete about when this ordering is needed at all, because a
great many tree edits never reach for it. Inserting a \emph{new} node is the
case where it is free. Set the node's parent pointer during the build phase,
while the node is still unpublished and no reader can reach it, and then make it
reachable with a single store to the child pointer in its parent. The node
becomes visible with its parent pointer \emph{already} correct, so there is one
publication rather than two, and a reader that descends into it and walks back up
finds the two edges agree with no ordering rule in play. This is ordinary \rcu
build-and-publish (\cref{sec:buildpublish}); it needs neither the reverse-order
rule nor a \pseudotxn.

The rule --- and the \fliplatch behind it --- earns its keep in the case
build-and-publish cannot reach: restructuring a node that is \emph{pinned}, its
address held by the application and so not free to be rebuilt hidden and swapped
in. A new node could carry its parent pointer into visibility because it was
assembled in private; a pinned node cannot, so its parent pointer has to be
re-pointed \emph{in place}, while readers can still reach it, and the forward and
back edges genuinely move under live readers. That is the re-parent above, and it
is where an atomic switch of both edges --- or, if the edit is split
across commits, the reverse-order rule --- stops being optional.

Which of the two applies is decided by the structure's traversal orders. The
reverse-order rule needs a \emph{single} fixed order to reverse; a cyclic
structure has none. A doubly-linked list read forward and read backward resolves
its two edges in opposite relative orders, so no commit order satisfies both
directions --- the bidirectional list of \cref{sec:wrappers}, and any tree walked
both up and down, must therefore switch all their reader-visible edges
atomically. The atomic commit is the general answer and the \fliplatch's
reason for being; the reverse-order rule is the cheaper option available only
where one order is fixed, as in the descend-then-ascend reader above.

The limits should be equally clear. This is an embedder's discipline, not an
facility guarantee: the facility cannot know a structure's traversal order, so
nothing checks that the commits were ordered against it. It applies only where a
read order actually is fixed; a reader that compares two independently-reached
slots, or re-reads a pointer it cached, has no chain to carry the order and is
back to owing itself an acquire. And it buys coherence along the traversal, not
a snapshot: the reader still cannot ask what the whole structure looked like at
one instant. What it can do --- and what a very large class of traversals needs
--- is never be told that time ran backwards.

That is what converts \cref{prop:monotone} from a consolation into a tool, and
it is the honest reply to the reviewer who asks why anyone would accept a model
without opacity.

\subsection{Why this weakening}
\label{sec:whyweakening}

\rcu exists because readers dominate.\footnote{Read dominance is the typical
case, not a rule: \rcu also serves update-heavy uses --- as a deferred reclaimer
that sidesteps the ABA problem, or where the read path must meet a real-time
deadline whatever the write rate. A \pseudotxn is shaped for the read-dominated
case regardless.} Its entire contract is that a reader pays
nothing: no lock, no barrier, no write to shared state. Every design that adds
multi-pointer atomicity on top of \rcu must decide what to charge that reader,
and the answer is what separates them (\cref{tab:comparison}).

% ---------------------------------------------------------------------------
% GUARDRAIL (do not delete in an editing pass).
%
% The "only once you mutate in place" qualifier is load-bearing, not hedging.
% An earlier draft asserted flatly that "isolation is paid for on the read side"
% and argued it from "a snapshot is a claim about what a reader observes, so
% only the reader can establish it".
%
% That premise is FALSE, and this paper of all papers cannot afford it: Sec. 1
% already names copy-on-write as one of RCU's three remedies. A writer that
% path-copies and swaps a root gives readers a fully isolated immutable
% snapshot at ZERO read-side cost -- the reader establishes nothing. Persistent
% data structures are the existence proof, and RCU's own build/publish/reclaim
% is the same trick in the small. The unqualified claim contradicts our own
% Sec. 1 and a reviewer who knows persistent structures spots it immediately.
%
% The true claim is narrower and still sufficient: isolation is a read-side
% charge for designs that mutate IN PLACE, which is the regime every row of
% tab:comparison occupies, ours included. Do not "tighten" this back into the
% unqualified form -- it reads punchier and it is wrong.
% ---------------------------------------------------------------------------
\paragraph{Isolation is a read-side cost only once you mutate in place.} The
claim needs its precondition stated, because without it the claim is false and
this paper has already said why. There is a way to give readers a snapshot for
nothing, and \rcu uses it: copy. A writer that builds a new version where
readers cannot reach it and then swaps one pointer hands every reader an
immutable structure. The reader establishes nothing, validates nothing, and
still sees a coherent view of everything reachable from the root it loaded.
Persistent data structures make this the entire design, and \rcu's
build--publish--reclaim discipline (\cref{sec:buildpublish}) is the same idea in
the small --- which is exactly why an \rcu reader gets a coherent view of a
\emph{node} for free.

What copying does not do is scale to the update, and that is where this paper
began (\cref{sec:onestore}). When a change spans several words, copying a region
large enough to make it one pointer swap means paying for that region --- in
allocation, in footprint, and in reclamation --- on every update. The
alternative is to mutate in place.

Mutating in place is what ends the free lunch, and it is the regime every design
in \cref{tab:comparison} occupies, this one included. A reader can now arrive
while the structure is between states. Any property it must observe about a
mutation in flight has become a claim about what \emph{that reader} sees, and
only the reader can establish it.
The established designs each make it reachable from the reader: an object
header, a version chain, a logged read set, a per-element field
(\cref{sec:costcomparison}).

So there are two ways to spare the reader, and this design takes the second. It
declines to copy, because the cost of copying is what made the operation
expensive in the first place. And it declines to hand readers isolation, because
in a structure mutated in place isolation is a charge on the reader --- and the
reader is the one party \rcu exists to keep free.

\Cref{sec:costcomparison} sets out what each of those alternatives charges the
reader, and why this design's read-side test instead rides a value the reader
already holds. Here we turn to what it charges the writer.

\paragraph{The cost moves to the write side.} This is not a free lunch, and the
evaluation must not pretend otherwise. Where existence collapses an element's
entire membership into a single object it can flip once, a \pseudotxn names
individual slots, so its commit is as wide as the number of edges it changes. A
structure whose logical update touches many pointers --- a tower, a wide node
--- therefore commits many records where existence commits one. Commit width is
the scarce resource, and it is the axis on which this design is at its worst.
% -----------------------------------------------------------------------
% NOTE (drafting; retained deliberately -- see README framing rule).
% Say this plainly, and let the evaluation confirm it rather than apologize
% for it: the existence-vs-txn gap on a tower-structured workload should
% WIDEN with tower height, because width grows and existence's does not.
% That is a prediction the cost model makes, and it is testable with the
% facility's own commit-cost introspection. A validated cost model is a
% result; an unexplained loss is an embarrassment. They are the same
% numbers.
% -----------------------------------------------------------------------

\paragraph{The weakening is targeted, not residual.} We did not fail to provide
isolation; we declined to charge for it. The structures \rcu serves are
read-dominated by construction, so paying permanently, in every element, for a
snapshot those readers never requested spends the one resource they cannot
spare. A \pseudotxn removes exactly that charge and relocates the cost onto the
writers, who are by hypothesis the minority.

One qualification belongs here rather than in a footnote, because the argument
above is otherwise too flattering. What RLU buys with its charge is strictly
more than what we provide: a coherent snapshot, where we give only monotone
resolution (\cref{prop:monotone}). "We declined to charge for it" is therefore
an honest description of the trade only if it travels with what the charge
bought. A reader who needs a snapshot cannot take this facility and economize ---
they need a different facility.

What can be said across that boundary is what the stronger guarantee costs the
\emph{reader}, and \cref{tab:readclass} says it: 27--50\% against RLU, and
14--21\% against existence, whose per-element guarantee is likewise stronger
than ours. Those figures are sound in the only direction that matters here,
because they run in our favor and against the grain of the trade --- the other
schemes charge the reader more \emph{and} deliver more, so the gap is a price,
not a defect. What would not be sound is ranking write throughput across the
same boundary, or reading these numbers as a claim that the facilities are
interchangeable; they are not (\cref{sec:relrcu}).

% ---------------------------------------------------------------------------
% GUARDRAIL (do not delete in an editing pass).
%
% The paragraph above is the honest asterisk on Sec. 3.7's whole argument, and
% it must stay IN THE TEXT rather than become a footnote. RLU's tax buys
% strictly MORE than ours does -- a coherent snapshot, where we give "never
% new-then-old" -- so a comparison against RLU crosses guarantee classes and
% must say so.
%
% REVISED 2026-07-26, and the distinction is the point: crossing is sound on
% the READER AXIS and only there. tab:readclass reports what the stronger
% guarantee costs the reader, and those numbers run AGAINST the grain of the
% trade -- the other schemes charge readers more and deliver more -- so they
% cannot flatter us by construction. Ranking WRITE throughput across the same
% boundary would not be sound, and neither would any reading that treats the
% facilities as interchangeable; both are ruled out in the text above.
% Do NOT let an editing pass restore a blanket "no comparison against RLU":
% that was the earlier scope, and it left Sec. 3.7's central trade unquantified.
% Sec. 9 states the same asterisk for the measurement specifically; both are
% needed, because Sec. 3.7 is where a reader forms the impression that the
% trade was free.
% ---------------------------------------------------------------------------

% ===========================================================================
\section{The Single-Writer Flip-Latch}
\label{sec:fliplatch}

\subsection{The intermediate-state problem}
\label{sec:intermediate}

A structural \rcu mutation often has to re-point many slots at once. Re-parenting
a set of live nodes must rewrite each node's back-pointer; a reader walking
several of them must not see a mix of old and new targets. Re-pointing the slots
one at a time cannot offer that: each store is individually atomic, but between
them the set is part-old and part-new, and a reader traversing the set observes
exactly that.

The conventional remedy is to drain the readers --- a grace period around the
update, so that no reader spans the inconsistent window. It is correct, and it
is expensive: the writer waits out every reader in the system to publish a
change none of them were likely looking at.

\subsection{Mechanism: one level of indirection}
\label{sec:swmechanism}

The flip-latch removes the intermediate states rather than hiding them behind a
drain, and it does so with one level of indirection
(\cref{fig:fliplatch}).

Each slot in the set is made to hold a \emph{tagged} pointer to a small
\emph{proxy} rather than the target directly. A proxy carries
$\{\mathit{old},\mathit{new}\}$ and a pointer to a shared \emph{flip group}
carrying one boolean \emph{selector} word. Resolving a proxy returns
$\mathit{old}$ while the selector is $0$, and $\mathit{new}$ once it is $1$.

Everything follows from one property: \textbf{every proxy in a group reads the
same selector.} A single store flipping that selector therefore switches the
whole set at once, and at every instant every proxy resolves consistently ---
all old, or all new. There is no instant at which the set is part-old and
part-new, because there is only one word that decides, and one store that
changes it.

% The flip-latch, installed: N slots -> N proxies -> ONE shared selector.
% This is the paper's essential diagram; everything in Sec. 4 refers to it.
%
% Drawn in TikZ (not an imported image) so the figure ships as source in the
% arXiv submission and needs no external toolchain to rebuild.
%
% The picture must make exactly one thing obvious: every proxy reads the SAME
% selector word, so the single store that flips it switches the whole set. If a
% reader of the figure can see that, the mechanism is explained.
\begin{figure}[t]
\centering
\begin{tikzpicture}[
  font=\small,
  slot/.style   = {draw, minimum width=15mm, minimum height=6mm, fill=gray!8},
  proxy/.style  = {draw, minimum width=26mm, minimum height=9mm, fill=blue!6},
  target/.style = {draw, rounded corners=2pt, minimum width=11mm, minimum height=5.5mm},
  old/.style    = {target, fill=gray!15},
  new/.style    = {target, fill=green!12},
  sel/.style    = {draw, very thick, minimum width=17mm, minimum height=8mm, fill=orange!18},
  arr/.style    = {-latex, thin},
  dep/.style    = {-latex, thin, densely dotted, gray!70!black},
]

% ---- slots (the words the reader loads) -------------------------------------
\node[slot] (s1) at (0, 1.6)  {\code{s$_1$}};
\node[slot] (s2) at (0, 0)    {\code{s$_2$}};
\node[slot] (s3) at (0, -1.6) {\code{s$_3$}};
\node[above=1.5mm of s1, align=center, font=\footnotesize\itshape] {slots\\(tagged)};

% ---- proxies ---------------------------------------------------------------
\node[proxy] (p1) at (3.6, 1.6)  {$\{o_1,\, n_1\}$};
\node[proxy] (p2) at (3.6, 0)    {$\{o_2,\, n_2\}$};
\node[proxy] (p3) at (3.6, -1.6) {$\{o_3,\, n_3\}$};
\node[above=1.5mm of p1, align=center, font=\footnotesize\itshape] {proxies};

\draw[arr] (s1) -- (p1);
\draw[arr] (s2) -- (p2);
\draw[arr] (s3) -- (p3);

% ---- the shared selector: the whole point ----------------------------------
\node[sel] (sel) at (7.9, 0) {\code{sel}: $0 \to 1$};
\node[above=1.5mm of sel, align=center, font=\footnotesize\itshape] {flip group\\(one word)};

\draw[dep] (p1.east) -- (sel.north west);
\draw[dep] (p2.east) -- (sel.west);
\draw[dep] (p3.east) -- (sel.south west);

% ---- resolution: the WHOLE old set, or the WHOLE new set ------------------
% Not one row expanded and the other two implied -- that misses the point of
% the figure, which is that the selector moves the entire set together. Two
% arrows, two groups; the reader sees three slots resolving as one.
\node[old] (o1) at (11.3,  2.45) {$o_1$};
\node[old] (o2) at (11.3,  1.75) {$o_2$};
\node[old] (o3) at (11.3,  1.05) {$o_3$};
\node[draw, densely dashed, gray!60, inner sep=2.2mm, fit=(o1)(o2)(o3)] (oldset) {};

\node[new] (n1) at (11.3, -1.05) {$n_1$};
\node[new] (n2) at (11.3, -1.75) {$n_2$};
\node[new] (n3) at (11.3, -2.45) {$n_3$};
\node[draw, densely dashed, green!45!black, inner sep=2.2mm, fit=(n1)(n2)(n3)] (newset) {};

% The selector value labels the GROUP, not the arrow. On the path a label is
% crossed by its own line; beside the path it collides with the "flip group"
% caption. On the group it is unambiguous (the arrow points at the group the
% condition selects) and it cannot collide with anything.
\draw[arr, gray!70!black] (sel.east) -- (oldset.west);
\draw[arr, thick]         (sel.east) -- (newset.west);
\node[font=\scriptsize, anchor=south] at (oldset.north) {\code{sel}=0};
\node[font=\scriptsize, anchor=south] at (newset.north) {\code{sel}=1};

\node[align=center, font=\footnotesize\itshape, anchor=north] at (11.3, -3.05)
  {one word decides\\the whole set};

% ---- the flip --------------------------------------------------------------
\draw[-latex, very thick, orange!70!black]
  (sel.south) ++(0,-0.15) -- ++(0,-0.75)
  node[below, align=center, font=\footnotesize\bfseries, text=black]
  {the commit:\\one release store};

\end{tikzpicture}
\caption{The flip-latch, installed, with three slots standing for $N$. Each slot
holds a \emph{tagged} pointer to a proxy rather than to its target; each proxy
carries $\{\mathit{old},\mathit{new}\}$ and reads one shared selector word.
While \code{sel} is $0$ every proxy in the group resolves to its old target, so
installation is invisible to readers. The commit is a single release store,
$0 \to 1$, after which every proxy resolves to its new target. The two dashed
groups are what one store chooses between: there is no instant at which the set
is part-old and part-new, and no per-slot work at commit time.}
\label{fig:fliplatch}
\end{figure}


This is the commit step of \cref{def:pseudotxn}, realized. The record set is the
set of proxies to be installed; it is frozen when installation begins, because
from that point the writer is placing proxies for exactly those records and no
other may join. Installation itself publishes nothing --- the selector is still
$0$, so every proxy resolves to old --- and the commit step is the lone selector
store that follows. That single store, $0 \to 1$ and never written back, is also
why \cref{prop:monotone} holds.

\subsection{Why exclusion is required}
\label{sec:exclusion}

The flip-latch is legal \emph{only} under writer mutual exclusion. This is a
precondition of the mechanism, not a caveat on it.

The facility has no install-time CAS, no conflict detection, and no abort. Those
are precisely the parts that would let two writers discover each other, and they
are absent by design --- their absence is what makes this facility cheaper than the
multi-writer one. The consequence is unforgiving: two writers racing on one slot
simply corrupt it, and a default build says nothing.
% -----------------------------------------------------------------------
% NOTE (drafting; retained deliberately -- see README framing rule).
% Do not soften this sentence. It is the whole contract, and an embedder who
% misreads it writes a corruption that no test in a default build will
% surface.
% -----------------------------------------------------------------------

Because the failure is silent, the facility ships a runtime validator
(\code{-DURCU\_TXN\_SW\_EXCL\_VALIDATE}) that aborts the process on a violation.
It is off by default and costs nothing when off --- the handle does not even
carry the owner field. It claims the two things that exist to be claimed: the
\emph{handle}, whose owner is the thread that initialized it (catching a
\pseudotxn handed between threads mid-flight), and the \emph{slot}, which is
what two racing writers actually share and so is where the violation is really
visible. In a correct single-updater program a slot being recorded or parked
cannot already hold a parked proxy.
% -----------------------------------------------------------------------
% NOTE (drafting; retained deliberately -- see README framing rule).
% This is worth stating as a design lesson rather than an implementation
% note: the contract is unenforceable at zero cost, so the facility makes the
% check opt-in and free-when-off rather than pretending the contract is
% safe. Ties to the series-wide honesty point: memory safety is
% unconditional, atomicity is conditional on embedder contracts. Cross-
% reference Sec. 7.
% -----------------------------------------------------------------------

\subsection{Lifecycle: build, install, commit, settle}
\label{sec:lifecycle}

The mutation is driven by the \emph{embedder}: the code that builds a data
structure on this facility and performs its updates. The embedder owns the whole
lifecycle---all four steps below---and the obligations of
\cref{sec:embedderowes}; "builder" would name only the first step.
\begin{enumerate}
\item \textbf{Build} the new structure where readers cannot reach it.
\item \textbf{Install.} The record set is frozen at this point: these are the
      slots that will change, and no further record may join. Initialize the
      group, then for each slot initialize a proxy with its
      $\{\mathit{old},\mathit{new}\}$ targets and point the slot at the tagged
      proxy. The selector is $0$, so this is \emph{transparent to readers}: they
      still resolve to old. Installation is not a partial commit; it is not
      observable at all.
\item \textbf{Commit.} One release store, $0 \to 1$. Every proxy now resolves to
      new, atomically. This is the linearization point
      (\cref{prop:lp}).
\item \textbf{Settle.} Rewrite each slot from the tagged proxy back to the
      direct new target. This is idempotent for readers, since the proxy already
      resolves to new --- so settle changes no observable value and is not a
      linearization point either. Then \code{call\_rcu()} the proxies and the
      group.
\end{enumerate}

The shape is exactly build~$\rightarrow$~publish~$\rightarrow$~reclaim
(\cref{sec:background}), with the commit as the single publish. Steps 2 and 4
bracket it symmetrically: install makes the indirection reader-visible without
changing any value, and settle removes it without changing any value.

The claim that only one of these four steps is observable is easy to assert and
worth checking. \Cref{tab:lifecycle} tracks a single transacted slot $s$ through
the lifecycle, alongside what a concurrent reader resolving $s$ actually
obtains.

\begin{table}[tb]
\centering
\caption{One transacted slot $s$ through the lifecycle. The reader's column
changes exactly once. Install and settle both rewrite $s$ --- they are stores to
a word readers are concurrently loading --- yet neither changes what resolving
$s$ yields, which is why the commit is the only linearization point.}
\label{tab:lifecycle}
\begin{tabular}{@{}llll@{}}
\toprule
Step & $s$ holds & Reader resolving $s$ obtains & Linearization point? \\
\midrule
(before)  & $o$, directly            & $o$ & --- \\
Build     & $o$, directly            & $o$ & no: new nodes unreachable \\
Install   & tagged proxy, \code{sel}$=0$ & $o$ & no: resolved value unchanged \\
Commit    & tagged proxy, \code{sel}$=1$ & $n$ & \textbf{yes} \\
Settle    & $n$, directly            & $n$ & no: resolved value unchanged \\
(after)   & $n$, directly            & $n$ & --- \\
\bottomrule
\end{tabular}
\end{table}

The third column is the mechanism. It reads $o, o, o, n, n, n$: one transition,
at the commit, no matter how many slots the set contains or how long the writer
takes over the other three steps. A reader cannot tell installed-but-uncommitted
from untouched, and cannot tell settled from committed-but-unsettled. Only the
selector store is visible, and it is one store.

\paragraph{The wrinkle, which is the contract in disguise.} Across the install
window a proxy resolves to the old target \emph{recorded at install time} rather
than to whatever the slot happens to hold. Under exclusion those are the same
thing --- no one else writes $s$, so the recorded $o$ stays the true $o$ --- and
the row for Install above is exact.

Remove exclusion and that row becomes a lie in the most inconvenient way
available. A second writer's store to $s$ during the window is not detected and
not lost-but-noticed: the proxy goes on publishing the $o$ it recorded, so
readers resolve to a value that is no longer what the slot was set to, and the
subsequent settle overwrites the intruder's store with $n$. Nothing traps, and
readers are handed a stale value \emph{as though it were current}. This is the
same fact as \cref{sec:exclusion} viewed from the reader's side, and it is worth
seeing from both: the \fliplatch's cheapness comes from having no install-time
check. A facility that checked at install --- a compare-and-swap on the slot,
paid on every commit --- would have caught the second writer and refused the
install; this one spends nothing there and asks the embedder for exclusion
instead.

\subsection{Progress}
\label{sec:swprogress}

Under exclusion there is nothing to contend for, and the progress story is
correspondingly short: a commit cannot fail. There is no abort, because there is
no contention to abort on. A committed \pseudotxn is the writer's to complete;
nothing a concurrent party does can turn it back.
% -----------------------------------------------------------------------
% NOTE (drafting; retained deliberately -- see README framing rule).
% An earlier draft explained here why the API carries an ABORT status value the
% single-updater never returns (it is shared with the concurrent front-end so an
% embedder migrates without rewriting its commit handling). Removed: it explains
% a P1 feature by a P2 mechanism the reader has not met, and it is redundant --
% the migration point already lives in Sec. 1's Scope paragraph, which is where
% the series relationship belongs. Do not reintroduce a concurrent-facility
% explanation in this subsection; keep the progress story self-contained.
% -----------------------------------------------------------------------

\subsection{What the embedder owes}
\label{sec:embedderowes}

The facility is deliberately incomplete: it supplies the latch, and leaves three
obligations with the structure that embeds it.

\begin{itemize}
\item \textbf{Tagging.} The embedder chooses how a proxy pointer is marked so
      its readers recognize one --- a spare low or high bit, or a reserved type
      code in an already-tagged pointer --- and routes resolution through the
      facility's accessor. \Cref{sec:tags} takes this up: it is where the facility's
      zero-added-state claim is made concrete.
\item \textbf{Allocation.} Proxies and the group are the embedder's to allocate,
      with whatever alignment its tagging scheme demands --- often as a single
      backing block with one \code{rcu\_head}.
\item \textbf{Reclamation.} After settle has unpublished every proxy, the
      embedder \code{call\_rcu()}s them.
\end{itemize}

% -----------------------------------------------------------------------
% NOTE (drafting; retained deliberately -- see README framing rule).
% Frame this as a design position, not a shortfall: the facility cannot know
% the embedder's pointer encoding, so a general facility that OWNED the
% tagging would have to add a field -- which is the very cost this design
% exists to avoid. The obligations are the price of adding no per-element
% state, and that trade is the paper's thesis. Do not let this read as "and
% the user has to do the hard part".
% -----------------------------------------------------------------------

% ===========================================================================
\section{Slots and Tags}
\label{sec:tags}

\Cref{sec:whyweakening} claimed that this facility adds no per-element state and
costs the reader no additional load. This section substantiates that claim, and
states what it costs in return: the facility has nowhere of its own to put the
marker, so it must borrow bits from the embedder's pointer encoding.

\subsection{The tag contract}
\label{sec:tagcontract}

A parked slot must hold something a reader can recognize as a proxy and
distinguish from every legitimate value the slot could otherwise hold. The
facility takes those bits from the slot itself:

\begin{itemize}
\item a parked slot's value is \code{(record\_address | tag)};
\item the record is recovered as \code{(value \& \textasciitilde tag)};
\item a value is this record's proxy iff \code{(value \& tag) == tag}.
\end{itemize}

Which imposes a contract on the embedder, and it is unforgiving:
\textbf{every value the slot may legitimately hold must satisfy
\code{(value \& tag) != tag}.} A single value that happens to carry the tag
pattern is indistinguishable from a proxy, and a reader will follow it into a
record that does not exist.

A debug build (\code{URCU\_TXN\_SW\_EXCL\_VALIDATE}) does not leave this to
inspection: before parking a proxy on a single-writer slot, it asserts the slot
does not already read as parked. Its first purpose is to catch a second writer
on a slot the caller promised exclusive; but the same test trips on an
ill-chosen tag, since a legitimate value carrying the tag pattern reads as a
parked proxy --- so under a single writer, where no second writer exists, that
abort is the aliasing tag itself, surfaced.

Two further obligations follow. Every write to a transacted slot must go through
the facility: one raw store breaks atomicity for every participant, not only the
writer that made it. And every read must resolve through the facility's accessor,
since a bare load of a parked slot yields a tagged record address rather than a
value.

\subsection{Why the tag cannot be fixed}
\label{sec:pertag}

The obvious design is a single tag, chosen once by the facility --- bit~0, say,
free in any aligned pointer. It is also the design that
\cref{prop:crossstructure} rules out.

Bit~0 is free only in structures that have not already spent it. A structure
whose pointers encode a type or a mark in their low bits has not; for such a
structure, a bit-0 proxy \emph{aliases a real value}, and the contract of
\cref{sec:tagcontract} is violated by construction. Such a structure can still
be transacted --- it simply has to nominate a different pattern, one that no
value of its own can carry, perhaps a reserved low nibble rather than a single
bit.

That alone would argue for a per-embedder tag, which a compile-time constant
could supply. What forces the tag \emph{per record} is cross-structure
atomicity. One \pseudotxn may span structures with different encodings: the
whole point of \cref{prop:crossstructure} is that a single commit can publish a
node into one structure and splice it into another --- one tagging a single bit,
the other a reserved nibble because its own pointers already spend the low bits.
Those slots have different free bits, so \emph{one transaction needs two tags at
once}. A per-translation-unit macro cannot express that; the tag has to travel
with the record.

So each record carries its own \code{proxy\_tag}, supplied when the record is
added, and heterogeneous slots share one facility, each resolved through its own
tag.

Three things must be kept apart here: the tag's \emph{meaning}, the slot, and the
record. The meaning belongs to the slot's structure --- not to the record, and
not to the value the slot holds. A reader never reads the tag out of that value;
it knows which structure it is traversing, hence that slot's tag, statically ---
the typed accessor tests a compile-time constant, the generic load takes the tag
as a parameter. The record keeps its own copy of the tag only for the write
side: the generic commit machinery, blind to structure types, reads each
record's copy to park each proxy with the pattern its slot's readers will test
for --- which is what lets one commit carry a bit-0 list edge and a nibble-tagged
structural edge together, each record with its own copy.
% -----------------------------------------------------------------------
% NOTE (drafting; retained deliberately -- see README framing rule).
% This is the tightest argument in the paper: Property 4 (cross-structure
% atomicity) FORCES the per-record tag. It is not a convenience or a
% generalisation-for-its-own-sake -- a per-TU tag would make Property 4
% unimplementable for heterogeneous embedders. Lead with that, not with
% "flexibility". Do NOT use the fractal trie as the worked example here --
% it is deferred to a later paper and P1 must not depend on it. The hazard
% is fully statable in the abstract: "a structure whose pointers already
% encode a type in their low bits". The facility header names the trie's 0xF
% nibble as its motivating case; we can describe the shape without the
% consumer.
% -----------------------------------------------------------------------

\subsection{Storing the tag, not a tagged pointer}
\label{sec:tagnotptr}

A record stores the \emph{tag}, not a pre-tagged pointer, and the proxy is
formed from the record's current address at install time. That distinction looks
pedantic and is not: it is what makes a descriptor safe to grow.

Growth here means reallocation, and reallocation moves addresses. A descriptor
that fills up is enlarged the way any dynamic array is --- a larger block is
allocated, the records are copied into it, and the old block is freed --- so
every record lands at a \emph{new} address. Had each record cached a proxy value
computed from its old address, the copy would strand all of them at once: each
would carry a pointer into storage the records no longer occupy. Because the
record stores only the tag, the tag travels with the record's content across the
move untouched, and the installer re-derives the proxy from the record's final
address. Nothing is stranded, and growth needs no fixup pass.

\subsection{The alignment budget}
\label{sec:alignment}

The tag's width is bounded by how many low bits of a record address are
guaranteed clear. Records are 16-byte aligned, so an embedder has four bits to
choose from: enough for a single bit, or for a reserved nibble.\footnote{Storing
a tag in a pointer's spare low bits, as this facility does, rests on
implementation-defined casts between pointers and integers; a proposal to give
C++ a standard pointer-tagging facility is under review~\citep{wg21p3125}.}

The alignment is established deliberately rather than inherited. The descriptor
is allocated with \code{posix\_memalign()} to a 16-byte boundary, and the record
array's offset and stride are both multiples of 16, so every record address ---
including every inline record --- has its low four bits clear. This holds
\emph{portably}: relying on \code{malloc} would rely on an allocator's habits
rather than a guarantee, and \code{max\_align\_t} is only 8 bytes on some 32-bit
ABIs.
% -----------------------------------------------------------------------
% NOTE (drafting; retained deliberately -- see README framing rule).
% Small but worth one sentence: "we asserted the alignment we need rather
% than inheriting whatever the allocator happened to give us" is exactly the
% kind of detail an enabling disclosure needs, and exactly the kind a normal
% paper drops. Keep it.
% -----------------------------------------------------------------------

\subsection{What the tag buys}
\label{sec:tagbuys}

The tag is the whole of the facility's presence in a reader's world. There is no
header, no side table, no per-element field --- the marker occupies bits of a
word the structure already had and the reader already loads. That is why
\cref{tab:comparison}'s last row reads zero in the loads column, and it is why
the structure a
reader traverses in the steady state is byte-for-byte the one it would traverse
with no facility linked at all.

The price is that the facility cannot own the encoding. It does not know which
bits a given structure has left, so it cannot choose for the embedder ---
\cref{sec:embedderowes}'s obligations are the direct consequence. A facility that
owned the tagging would have to add a field to say what it owned, which is
precisely the cost this design exists to avoid.
% -----------------------------------------------------------------------
% NOTE (drafting; retained deliberately -- see README framing rule).
% Close the loop with Sec. 3.7 explicitly: the obligations are not a
% shortfall, they are the mechanism by which the cost is avoided. A facility
% that took the decision away from the embedder would have to pay for the
% privilege in exactly the coin (per-element state) the design refuses to
% spend.
% -----------------------------------------------------------------------

% ===========================================================================
\section{Pseudo-Transacted Data Structures}
\label{sec:wrappers}

\Cref{sec:onestore} established the gap this facility exists to close: a
doubly-linked list has two reader-visible chains, but \rcu can publish only one
of them atomically, so \code{rculist} maintains a \code{prev} chain no reader may
walk and ships no \rcu-safe reverse iterator. The \fliplatch moves both edges at
a single instant --- precisely what that gap needs. This section spends it on two
structures: a bidirectional list whose reverse walk readers can take, and a hash
list.

\subsection{Two edges, one flip}
\label{sec:swlist}

Every structural change to a doubly-linked list re-points exactly two
reader-visible edges, one forward and one backward:
\begin{center}
\begin{tabular}{@{}ll@{}}
insert $E$ between $P$ and $N$: & \code{P->next}: \edit{N}{E} \quad and \quad \code{N->prev}: \edit{P}{E} \\
remove $E$ between $P$ and $N$: & \code{P->next}: \edit{E}{N} \quad and \quad \code{N->prev}: \edit{E}{P} \\
\end{tabular}
\end{center}

Two edges is exactly the case \cref{sec:swmechanism} handles. Both slots
transiently hold tagged proxies sharing one selector word; one release store
flips the selector and switches both edges together. At every instant the
forward and backward chains describe the same list, so the reverse walk that
\code{rculist} cannot offer becomes one a reader may take --- within the
guarantee, and its scope, that \cref{sec:coherence} states next.

% Inserting node E between P and N of a bidirectional list, P <-> N  =>  P <-> E <-> N.
%
% Three rows, read as a BUFFER between updater and readers:
%   Reader (top)     -- the only row carrying the list itself: one coherent list
%                       before the commit, another after, both directions.
%   RCU pseudo-txn   -- the buffer, drawn as its lifecycle install|commit|settle,
%     (middle)          aligned under the before | commit | after columns.
%   Updater (bottom) -- the single semantic mutation it issues.
% Insert removes nothing, so there is no grace period and no reclaim: the edit
% is only the atomic publish. A record is written concisely as <slot: old->new>
% (see the legend); Fig.~\ref{fig:fliplatch} carries the full proxy/selector anatomy.
%
% Self-contained TikZ (no external toolchain), matching fig-fliplatch.tex.
\begin{figure}[t]
\centering
\begin{tikzpicture}[
  font=\small,
  lnode/.style   = {draw, rounded corners=2pt, minimum width=8mm, minimum height=5.5mm, fill=gray!8},
  nnode/.style   = {draw, rounded corners=2pt, minimum width=8mm, minimum height=5.5mm, fill=green!14},
  record/.style  = {draw, rounded corners=1pt, fill=blue!6, inner sep=3.5pt, minimum width=35mm, minimum height=6.5mm, font=\footnotesize},
  settled/.style = {draw, fill=gray!8, inner sep=3.5pt, font=\footnotesize},
  sel/.style     = {draw, very thick, minimum width=13mm, minimum height=8mm, fill=orange!18},
  rowlab/.style  = {align=center, font=\footnotesize\itshape, text=black!75},
  colhead/.style = {align=center, font=\footnotesize\itshape},
  phase/.style   = {font=\footnotesize\itshape, text=black!70},
  biarr/.style   = {latex-latex, semithick},
  arr/.style     = {-latex, thin},
  dep/.style     = {-latex, thin, densely dotted, gray!70!black},
  selbox/.style  = {draw, very thick, minimum width=33mm, minimum height=5mm, fill=orange!18, inner xsep=3pt, font=\footnotesize},
  trow/.style    = {anchor=west, inner sep=0pt, minimum height=6mm, font=\footnotesize},
  pick/.style    = {-latex, thin, gray!55},
  ptr/.style     = {-latex, semithick},
]

% ---- column headers (reader states) ----------------------------------------
\node[colhead] at (2.6, 6.35) {before commit};
\node[colhead] at (9.5, 6.35) {after commit};

% ---- neutral before | after boundary through the reader row ----------------
\draw[thin, gray!45, densely dashed] (7.3, 6.1) -- (7.3, 3.15);

% ---- separators & row labels ----------------------------------------------
\draw[thin, gray!45] (-1.9, 3.05) -- (11.9, 3.05);
\node[rowlab] at (-0.95, 4.70) {Reader\\(bidirectional)};
\node[rowlab] at (-0.95, 1.35) {Updater\\(RCU pseudo-\\transaction)};

% ===== ROW 1 : reader's view -- the only row with the full list stacks ======
\node[lnode] (r1P) at (2.6, 5.9) {\code{P}};
\node[lnode] (r1N) at (2.6, 3.5) {\code{N}};
\draw[biarr] (r1P) -- (r1N);

\node[lnode] (r2P) at (9.5, 5.9) {\code{P}};
\node[nnode] (r2E) at (9.5, 4.7) {\code{E}};
\node[lnode] (r2N) at (9.5, 3.5) {\code{N}};
\draw[biarr] (r2P) -- (r2E);
\draw[biarr] (r2E) -- (r2N);
\node[right=1mm of r2E, font=\scriptsize, text=green!45!black] {new};

% ===== ROW 2 : the transaction, phase by phase =============================
% build -- the descriptor as a table (slot | old | new); selector box above it
\node[selbox] (selB) at (2.50, 2.72) {selector:~$0$};
\node[trow, fill=blue!14] (hdr)  at (0.85, 2.10) {\makebox[13mm][c]{\textit{slot}}\makebox[10mm][c]{\textit{old}\,(0)}\makebox[10mm][c]{\textit{new}\,(1)}};
\node[trow, fill=blue!6]  (rec1) at (0.85, 1.50) {\makebox[13mm][c]{\code{P.next}}\makebox[10mm][c]{\code{N}}\makebox[10mm][c]{\code{E}}};
\node[trow, fill=blue!6]  (rec2) at (0.85, 0.90) {\makebox[13mm][c]{\code{N.prev}}\makebox[10mm][c]{\code{P}}\makebox[10mm][c]{\code{E}}};
\draw[thin, blue!45] (0.85, 0.6) rectangle (4.15, 2.4);
\draw[thin, blue!45] (2.15, 0.6) -- (2.15, 2.4);
\draw[thin, blue!45] (3.15, 0.6) -- (3.15, 2.4);
\draw[thin, blue!45] (0.85, 1.8) -- (4.15, 1.8);
% each record row reads the shared selector
\draw[pick] (rec1.west) to[out=180,in=180] (selB.west);
\draw[pick] (rec2.west) to[out=180,in=180] (selB.south west);

% install -- a copy of the list; its pointers redirected into the record rows.
% P/N sit at the right of the install span rather than its centre: at the centre
% the arrows were 6mm, shorter than the labels riding on them, so next/prev were
% squeezed against the table edge. prev is labelled below, so the two labels
% cannot stack on each other in the 6mm gap between the rows.
\node[lnode] (iP) at (5.75, 1.50) {\code{P}};
\node[lnode] (iN) at (5.75, 0.90) {\code{N}};
\draw[ptr] (iP.west) -- (rec1.east) node[midway, above, font=\scriptsize] {\code{next}};
\draw[ptr] (iN.west) -- (rec2.east) node[midway, below, font=\scriptsize] {\code{prev}};

% commit -- the same selector, flipped (aligned above the commit phase)
\node[selbox, minimum width=27mm] (selC) at (7.30, 2.72) {selector:~$0 \to 1$};

% settle -- copy-paste of the install layout: the pointers now point at E itself
\node[nnode] (sE) at (9.10, 1.20) {\code{E}};
\node[lnode] (sP) at (10.70, 1.50) {\code{P}};
\node[lnode] (sN) at (10.70, 0.90) {\code{N}};
\draw[ptr] (sP.west) -- (sE.north east) node[pos=0.55, above, font=\scriptsize] {\code{next}};
\draw[ptr] (sN.west) -- (sE.south east) node[pos=0.55, below, font=\scriptsize] {\code{prev}};

% ===== ROW 3 : updater -- the transaction lifecycle, as phase spans =========
\draw[thin] (0.2, 0.35) -- (11.4, 0.35);
\foreach \x in {0.2, 4.2, 6.2, 8.4, 11.4} { \draw[thin] (\x, 0.24) -- (\x, 0.46); }
\node[phase] at (2.50, 0.03) {build};
\node[phase] at (5.15, 0.03) {install};
\node[phase] at (7.30, 0.03) {commit};
\node[phase] at (9.90, 0.03) {settle};

% the commit store, rising from the commit phase into the flipped selector
\draw[-latex, very thick, orange!70!black] (7.3, 0.55) -- (selC.south);
% Raised clear of the install P/N, which now reach further right.
\node[left, font=\scriptsize, text=black] at (7.15, 2.15) {one store};

% ---- legend for the record notation ---------------------------------------
\node[anchor=west, align=left, font=\scriptsize, text=black!75] at (-1.9, -0.55)
  {Each record is a row \{\textit{slot},\,\textit{old},\,\textit{new}\}; the shared selector picks the column\\
   the readers see --- the \textit{old} value while it reads $0$, the \textit{new} value after the flip to $1$.};

\end{tikzpicture}
\caption{Inserting \code{E} between \code{P} and \code{N} of a bidirectional
list, \code{P}\,$\leftrightarrow$\,\code{N} becoming
\code{P}\,$\leftrightarrow$\,\code{E}\,$\leftrightarrow$\,\code{N}, through the
flip-latch. The reader's row carries the list itself: it traverses one whole
coherent list before the commit and another after --- in either direction, never
a half-linked node. The second row is the updater's RCU pseudo-transaction,
traced across its lifecycle. It \emph{build}s the descriptor: a record binding
each moving edge to its \{\textit{old},\,\textit{new}\} values, all reading one
shared selector. \emph{install} redirects the list's own \code{next}/\code{prev}
pointers into those records. \emph{commit} is the single release store that flips
the selector $0\!\to\!1$. \emph{settle} leaves the pointers holding the plain new
node \code{E}. The new node is built privately first, its own \code{next} and
\code{prev} set to \code{N} and \code{P} before it is reachable. Because nothing
is removed, no list node is reclaimed; only the transaction's own descriptor
records are, after a grace period (cf.\ Fig.~\ref{fig:fliplatch} for the full
proxy/selector anatomy).}
\label{fig:listinsert}
\end{figure}


% Removing node E from P <-> E <-> N of a bidirectional list, leaving P <-> N.
%
% Derived from fig-list-insert: same two rows (reader; updater's RCU
% pseudo-transaction as build|install|commit|settle), with the values swapped
% for removal (old = E, new = the bypassing neighbour) and -- the point of this
% figure -- a reclaim tail. Unlike insert, a node leaves the list: E is unlinked
% at the commit but a reader parked on it reads through its untouched links, so
% synchronize_rcu must drain that reader before free(E). That drain-and-free is
% ordinary RCU, AFTER the transaction, not part of it.
%
% The record table is drawn TWICE, before and after the commit, with the live
% column shaded. Drawn once it showed the records only in the state the commit
% is about to leave, so the one thing the commit does -- move which column the
% readers see -- had to be inferred. The second copy is also what displays the
% interval between the commit and the settle: same records, still parked, now
% resolving new. The slots are rewritten only at the settle to the right of it.
%
% Self-contained TikZ (no external toolchain), matching fig-fliplatch.tex.
\begin{figure}[t]
\centering
\begin{tikzpicture}[
  font=\small,
  lnode/.style   = {draw, rounded corners=2pt, minimum width=8mm, minimum height=5.5mm, fill=gray!8},
  ghost/.style   = {draw, rounded corners=2pt, minimum width=8mm, minimum height=5.5mm, fill=gray!8, dashed, text=gray!55, draw=gray!55},
  rowlab/.style  = {align=center, font=\footnotesize\itshape, text=black!75},
  colhead/.style = {align=center, font=\footnotesize\itshape},
  phase/.style   = {font=\footnotesize\itshape, text=black!70},
  biarr/.style   = {latex-latex, semithick},
  gedge/.style   = {latex-latex, thin, densely dashed, gray!55},
  selbox/.style  = {draw, very thick, minimum height=5mm, fill=orange!18, inner xsep=3pt, font=\footnotesize},
  trow/.style    = {anchor=west, inner sep=0pt, minimum height=6mm, font=\footnotesize},
  pick/.style    = {-latex, thin, gray!55},
  ptr/.style     = {-latex, semithick},
  live/.style    = {fill=orange!16},
]

% ---- column headers (reader states) ----------------------------------------
\node[colhead] at (1.95, 6.35) {before commit};
\node[colhead] at (8.50, 6.35) {after commit};
\node[colhead] at (12.85, 6.35) {after synchronize};

% ---- reader-row dividers: at the commit, and at the end of synchronize ------
\draw[thin, gray!45, densely dashed] (6.20, 6.1) -- (6.20, 3.15);
\draw[thin, gray!45, densely dashed] (12.20, 6.1) -- (12.20, 3.15);

% ---- separators & row labels ----------------------------------------------
\draw[thin, gray!45] (-2.0, 3.05) -- (13.55, 3.05);
\node[rowlab] at (-1.15, 4.70) {Reader\\(bidirectional)};
\node[rowlab] at (-1.15, 1.35) {Updater\\(RCU pseudo-\\transaction)};

% ===== ROW 1 : reader's view -- P<->E<->N before, P<->N after ================
\node[lnode] (r1P) at (1.95, 5.9) {\code{P}};
\node[lnode] (r1E) at (1.95, 4.7) {\code{E}};
\node[lnode] (r1N) at (1.95, 3.5) {\code{N}};
\draw[biarr] (r1P) -- (r1E);
\draw[biarr] (r1E) -- (r1N);

\node[lnode] (r2P) at (8.50, 5.9) {\code{P}};
\node[ghost] (r2E) at (8.50, 4.7) {\code{E}};
\node[lnode] (r2N) at (8.50, 3.5) {\code{N}};
\draw[gedge] (r2P) -- (r2E);
\draw[gedge] (r2E) -- (r2N);
\draw[biarr] (r2P.east) to[out=0, in=0] (r2N.east);
\node[left=1mm of r2E, font=\scriptsize, text=gray!55] {bypassed};

% after synchronize -- all readers drained; the list is a clean P <-> N
\node[lnode] (r3P) at (12.85, 5.9) {\code{P}};
\node[lnode] (r3N) at (12.85, 3.5) {\code{N}};
\draw[biarr] (r3P) -- (r3N);

% ===== ROW 2 : the transaction, phase by phase =============================
% build -- the records as a table (slot | old | new); "old" is the live column
\node[selbox, minimum width=27mm] (selB) at (1.95, 2.72) {selector:~$0$};
\fill[live] (1.70, 0.6) rectangle (2.70, 1.8);
\node[trow, fill=blue!14] (hdr)  at (0.60, 2.10) {\makebox[11mm][c]{\textit{slot}}\makebox[10mm][c]{\textit{old}\,(0)}\makebox[10mm][c]{\textit{new}\,(1)}};
\node[trow, fill=blue!6]  (rec1) at (0.60, 1.50) {\makebox[11mm][c]{\code{P.next}}\makebox[10mm][c]{\code{E}}\makebox[10mm][c]{\code{N}}};
\node[trow, fill=blue!6]  (rec2) at (0.60, 0.90) {\makebox[11mm][c]{\code{N.prev}}\makebox[10mm][c]{\code{E}}\makebox[10mm][c]{\code{P}}};
\draw[thin, blue!45] (0.60, 0.6) rectangle (3.70, 2.4);
\draw[thin, blue!45] (1.70, 0.6) -- (1.70, 2.4);
\draw[thin, blue!45] (2.70, 0.6) -- (2.70, 2.4);
\draw[thin, blue!45] (0.60, 1.8) -- (3.70, 1.8);
\draw[pick] (rec1.west) to[out=180,in=180] (selB.west);
\draw[pick] (rec2.west) to[out=180,in=180] (selB.south west);

% install -- the list's pointers redirected into the record rows
\node[lnode] (iP) at (4.45, 1.50) {\code{P}};
\node[lnode] (iN) at (4.45, 0.90) {\code{N}};
\draw[ptr] (iP.west) -- (rec1.east) node[midway, above, font=\scriptsize] {\code{next}};
\draw[ptr] (iN.west) -- (rec2.east) node[midway, above, font=\scriptsize] {\code{prev}};

% commit -- the same selector, flipped
\node[selbox, minimum width=20mm] (selC) at (5.60, 2.72) {selector:~$0 \to 1$};

% after the commit -- the SAME records, still parked, now resolving "new"
\node[selbox, minimum width=24mm] (selA) at (7.85, 2.72) {selector:~$1$};
\fill[live] (8.40, 0.6) rectangle (9.40, 1.8);
\node[trow, fill=blue!14] (ahdr)  at (6.30, 2.10) {\makebox[11mm][c]{\textit{slot}}\makebox[10mm][c]{\textit{old}\,(0)}\makebox[10mm][c]{\textit{new}\,(1)}};
\node[trow, fill=blue!6]  (arec1) at (6.30, 1.50) {\makebox[11mm][c]{\code{P.next}}\makebox[10mm][c]{\code{E}}\makebox[10mm][c]{\code{N}}};
\node[trow, fill=blue!6]  (arec2) at (6.30, 0.90) {\makebox[11mm][c]{\code{N.prev}}\makebox[10mm][c]{\code{E}}\makebox[10mm][c]{\code{P}}};
\draw[thin, blue!45] (6.30, 0.6) rectangle (9.40, 2.4);
\draw[thin, blue!45] (7.40, 0.6) -- (7.40, 2.4);
\draw[thin, blue!45] (8.40, 0.6) -- (8.40, 2.4);
\draw[thin, blue!45] (6.30, 1.8) -- (9.40, 1.8);

% settle -- the slots rewritten to their new values; E is bypassed (P <-> N)
\node[lnode] (sP) at (10.45, 1.85) {\code{P}};
\node[lnode] (sN) at (10.45, 0.65) {\code{N}};
\draw[ptr] (sP) to[bend left=40] node[right=1pt, font=\scriptsize] {\code{next}} (sN);
\draw[ptr] (sN) to[bend left=40] node[left=1pt,  font=\scriptsize] {\code{prev}} (sP);

% ---- reclaim tail: ordinary RCU, AFTER the transaction ---------------------
% synchronize: E is unlinked, but its own links still point at N and P, so a
% reader parked on E reads straight through -- drain it over the grace period
\node[ghost] (dE) at (11.55, 1.55) {\code{E}};
\fill[black!70] ([yshift=1.3mm]dE.north) circle (1.7pt);
\node[font=\scriptsize, above=2.2mm of dE, align=center] {parked\\reader};
% free: after the grace period, no reader on E -> reclaim E and the records
\node[ghost] (fE) at (12.85, 1.55) {\code{E}};
\draw[gray!55] (fE.north west) -- (fE.south east);
\draw[gray!55] (fE.south west) -- (fE.north east);
\node[font=\scriptsize, below=0.5mm of fE, align=center] {\code{free(E)}\\+ records};
\draw[-latex, thin, gray!55] (dE.east) -- (fE.west)
  node[midway, above, font=\scriptsize, text=gray!60] {drains};

% ===== the lifecycle ruler ==================================================
\draw[thin] (0.2, 0.35) -- (13.55, 0.35);
\foreach \x in {0.2, 4.1, 5.0, 6.2, 10.85, 12.20, 13.55} { \draw[thin] (\x, 0.24) -- (\x, 0.46); }
\node[phase] at (2.15, 0.03) {build};
\node[phase] at (4.35, 0.03) {install};
\node[phase] at (5.60, 0.03) {commit};
\node[phase] at (9.20, 0.03) {settle};
\node[phase] at (11.50, 0.03) {synchronize};
\node[phase] at (12.85, 0.03) {free};

% the commit store, rising from the commit phase into the flipped selector
\draw[-latex, very thick, orange!70!black] (5.60, 0.55) -- (selC.south);
\node[left, font=\scriptsize, text=black] at (5.45, 2.05) {one store};

% ---- legend for the record notation ---------------------------------------
\node[anchor=west, align=left, font=\scriptsize, text=black!75] at (-2.0, -0.72)
  {Each record is a row \{\textit{slot},\,\textit{old},\,\textit{new}\}; the shaded column is the one the readers see. The commit\\
   moves the shading and nothing else --- the records, and the slots parked on them, are unchanged by it.};

\end{tikzpicture}
\caption{Removing \code{E} from
\code{P}\,$\leftrightarrow$\,\code{E}\,$\leftrightarrow$\,\code{N} of a
bidirectional list, leaving \code{P}\,$\leftrightarrow$\,\code{N}, through the
\fliplatch. The reader's row carries the list: it traverses
\code{P}\,$\leftrightarrow$\,\code{E}\,$\leftrightarrow$\,\code{N} before the
commit and \code{P}\,$\leftrightarrow$\,\code{N} after --- in either direction,
never a half-linked node. The second row is the updater's \pseudotxn:
\emph{build} the records (each moving edge bound to its
\{\textit{old},\,\textit{new}\} values, here \textit{old}\,$=$\,\code{E} and
\textit{new} the bypassing neighbour, all reading one shared selector),
\emph{install} the list's pointers into those records, \emph{commit} the single
store that flips the selector $0\!\to\!1$, and \emph{settle} the slots to their
new values (\code{P.next}\,$=$\,\code{N}, \code{N.prev}\,$=$\,\code{P}). The
table appears twice because the commit changes nothing but which column is live:
between it and the settle the same records are still parked, now resolving
\textit{new}. Unlike insert, a node leaves the list. At the commit \code{E} is
unlinked, but its own \code{next} and \code{prev} still point at \code{N} and
\code{P}, so a reader already parked on \code{E} reads straight through them;
\code{synchronize\_rcu} waits out the grace period until that reader drains, and
only then is \code{free(E)} --- together with the transaction's records ---
safe. This drain-and-free is ordinary \rcu, after the transaction,
not part of it (cf.\ \cref{fig:fliplatch}).}
\label{fig:listremove}
\end{figure}


\subsection{The coherence guarantee, and its scope}
\label{sec:coherence}

\begin{property}[Mutual inverses]
\label{prop:coherence}
For any two nodes that are both members of the list at the instant they are
resolved, the links are mutual inverses: \code{a->next == b} implies
\code{b->prev == a}. No operation exposes a window in which one direction sees
the change and the other does not.
\end{property}

This does not freeze the list for a long-lived reader --- nothing here does, and
\cref{sec:nonguarantees} said so. A reader may still observe an operation that
commits during its walk. What is guaranteed is that whatever it observes is
internally coherent in both directions \emph{at each step}.

``At each step'' is the precise scope, and it is \cref{sec:monotone}'s
dependency-chaining caveat in concrete form. Coherence is a property of a single
resolution; carrying it across two resolutions needs an ordering edge between
them. A dependency-chained walk --- next-then-next, or a \code{prev} hop taken
\emph{from} the node the previous load returned --- has one on every supported
architecture, as does any build using the default C11 dereference. A reader
comparing two independently-reached slots, a cached pointer or a \code{prev} and
a \code{next} read from nodes it did not chain between, must supply its own
acquire.

\paragraph{Across updates, not a snapshot.} \Cref{prop:coherence} is per-update:
one commit switches its own two edges together, so no reader catches a single
edit half-applied. Whether two edits are \emph{separate} commits, though, is the
embedder's choice, and the facility's reach is that they need not be --- a single
commit can carry both (\cref{prop:crossstructure}). Removing \code{A} and \code{D}
from \code{A}\,$\leftrightarrow$\,\code{B}\,$\leftrightarrow$\,\code{C}\,$\leftrightarrow$\,\code{D}
in one \pseudotxn, their edge re-points sharing one selector, is atomic to every
reader: it sees \code{A}\,$\leftrightarrow$\,\code{B}\,$\leftrightarrow$\,\code{C}\,$\leftrightarrow$\,\code{D}
or \code{B}\,$\leftrightarrow$\,\code{C}, and no state between. That composition is
what plain \rcu cannot offer --- it publishes one edit at a time and must sequence
the two, so a reader can straddle them.

Split the removals into separate commits and the straddle returns: a forward walk
can hold \code{A} after \code{D} has left the tail, reading
\code{A}\,$\leftrightarrow$\,\code{B}\,$\leftrightarrow$\,\code{C}, a list that
existed at no instant; and because a reverse walk resolves the edges in the
opposite order, no commit order spares both directions --- reversing the two
removals only moves the phantom to
\code{B}\,$\leftrightarrow$\,\code{C}\,$\leftrightarrow$\,\code{D}.
\Cref{prop:monotone} and the publish-order rule (\cref{sec:publishorder}) keep a
single edit's edges ordered; neither makes two separate commits atomic. Where the
edits genuinely must be separate --- decided by different threads, or far apart in
time --- coherence across them is the embedder's to add: a grace period to drain
readers first; where one direction suffices, issuing them in that reader's reverse
order; or a scheme layered on the value, a generation counter or a double
traversal that retries on mismatch, which later papers take up.
% -----------------------------------------------------------------------
% NOTE (drafting; retained deliberately -- see README framing rule).
% Do not bury this. It is the same precondition as Property 5's scope
% paragraph, and stating it twice -- abstractly in Sec. 3.5, concretely here
% -- is correct rather than repetitive, because this is where an embedder
% will actually trip over it.
% -----------------------------------------------------------------------

\subsection{The hash list: one transacted edge}
\label{sec:swhlist}

The hash-bucket sibling is shaped like the kernel's hlist: the head is a single
8-byte pointer, so a bucket table costs half what sentinel-node heads cost.

Its backward link is the elegant part. A node's \code{pprev} is a pointer to the
\emph{slot that names it} --- \code{\&prev->next}, or \code{\&head->first} ---
rather than a pointer to the previous node. So the head's single pointer is just
another \code{next} slot, and \textbf{no operation carries a head special
case}: first-position insert and delete are the general code path with a
different slot address.

Readers walk forward only, so a structural operation has exactly one
reader-visible edge --- the \code{next} slot it re-points. That edge is
transacted; \code{pprev} is not. No reader ever loads \code{pprev}, and a single
updater has no peer to race it, so the backlink needs neither atomicity nor a
proxy: it is maintained by a plain store, written eagerly in the prepare step.
An hlist mutation therefore records \emph{one} edge and commits it with a lone
release store --- and for a stand-alone insert or delete there is no proxy to
allocate, no group, and no deferred reclamation, so it costs about what a bare
\code{rcu\_assign\_pointer} does.

The plain backlink is not a corner that composition pays for later. When two
edits land on one bucket, the second reads its neighbour's \code{pprev} directly
and sees the first edit's store, because under one updater a plain store is
immediately visible to that same updater --- so adjacent deletes compose without
the backlink ever entering the \pseudotxn. \emph{Deferring} the backlink store
and then reading it raw would be the bug, the hlist form of the adjacent-delete
trap of \cref{sec:swcomposition}: a later edit would re-point a slot inside an
already-unlinked node. Eager-and-raw agree on every committed outcome;
deferred-and-raw do not.

Why transact the forward edge at all, when one release store would publish it?
For the same reason the bidir list does: composability. Through the
\code{*\_prepare} forms that single edge folds into a larger cross-structure flip
--- \cref{prop:crossstructure} put to work --- so an operation that needs no
atomicity of its own is still expressed as a \pseudotxn, and pays almost nothing
for the option when it is used alone.

\subsection{Composing edits: read-your-own-writes}
\label{sec:swcomposition}

Folding several edits into one commit raises a hazard that has nothing to do
with concurrency, and everything to do with a writer reading state its own
earlier edit has already changed. Consider deleting two adjacent nodes,
$P \rightarrow E_1 \rightarrow E_2 \rightarrow N$, in one \pseudotxn. Removing
$E_1$ records \code{P->next}: \edit{E_1}{E_2} and \code{E2->prev}:
\edit{E_1}{P}. Removing $E_2$ must now record \code{P->next} again --- but only
if it \emph{sees} that \code{P->next} is already pending. If it re-reads the
live structure instead, it believes $E_2$'s predecessor is still $E_1$, records
\code{E1->next}: \edit{E_2}{N} on a node that is about to vanish, and the commit
publishes a list pointing through deleted nodes. Nothing here is a race; one
writer, editing in program order, has torn its own commit.

Neither the hazard nor its answer is ours. Every deferred-update transactional
memory meets it --- writes are buffered until commit, so a later read in the
same transaction would see stale state unless something intervenes --- and the
standard fix is that a transactional read consults the transaction's own write
set first and returns the buffered value when it finds one; object-based designs
reach the same end by giving the transaction a private working copy to read
through \citep{shavit1995stm, herlihy2003dstm, fraser2004lockfreedom,
dice2006tl2}. \emph{Read-your-own-writes} is that answer, adopted here unchanged
and claimed nowhere in \cref{sec:relclaims}. What is worth stating is where it
lands in a facility whose reads are not instrumented, which is the rest of this
section.

Two accessors form the pair: a load that returns a slot's pending new value if
this \pseudotxn has already recorded one and its committed value otherwise, and
a recording store that \emph{fuses} a second edit of a slot onto the existing
record --- keeping the committed old, advancing the new --- rather than
appending a duplicate. A transaction holds at most one record per slot. Keeping
the committed old is a constraint a redo log does not have and the \fliplatch
does: the record's old value is what a reader resolves the proxy to before the
commit, so it must stay the value the slot would have held all along, however
many times the writer revises the new one. The wrappers read every neighbour
link through the load and record through the fusing store, so the second delete
above reads the \emph{pending} \code{P->next}, chains onto that record, and the
commit unlinks both victims and leaves $P \leftrightarrow N$ coherent.

Two obligations remain on the caller, and the question of what the first of
them costs sits between them.

\emph{The search must go through the transactional load, or precede the edits.}
A traversal is nothing but a sequence of loads, so a caller that walks through
\code{urcu\_txn\_sw\_load} sees its own pending edits at every step, and may
interleave walking and editing freely: this is the same reconciliation that
serves the edits, applied to the steps between them. What it may not do is walk
through the ordinary \rcu accessors while staging edits into the same
\pseudotxn{} --- those resolve the committed structure, not the pending one, so an
edit can be handed a node an earlier edit has already displaced. The wrappers
take the cheaper route where it is available: the caller's initial lookup runs
on the ordinary accessors, before the first record is staged, and only the
neighbour links the edits touch are read transactionally. No shape of update is ruled out by this. The
choice is which accessor the search uses, and it is the caller's to make.

\emph{Walking transactionally costs a find per step, and the find is cheap.}
Every transactional load asks the same question the recording store asks ---
\emph{is this slot already pending?} --- and the phrasing invites a picture of
searching an array of proxy records, which would indeed not be fast. The search
is not that. The write set is private to the \pseudotxn and to the thread that
owns it: no atomics, no lock, no line another thread is touching. For a typical
one, a two-edge list operation, it is a handful of records,
where a linear scan beats anything cleverer, and the facility does exactly that
scan. Only when a write set grows long enough for the scan to dominate does it
arm a Bloom filter, which answers the find's overwhelmingly common outcome on a
traversal --- a miss, this slot is not one I have edited --- in constant time.
A false positive costs the scan the filter was meant to save and nothing more;
there are no false negatives, so no pending edit can be walked past. The filter
is built lazily, on the first lookup
that would scan a long write set rather than on the record that crosses the
threshold, so a \pseudotxn that stays small never pays for one, and a long
walk over a large write set pays a hash rather than a scan.

None of this asks the writer to guard its traversal against anyone else. Under
exclusion (\cref{sec:exclusion}) no other writer can change a slot, so a
transactional walk has only this writer's own pending edits to reconcile;
there is no dependency a peer could move underneath it, and no reason to
transact a slot merely because the walk stepped through it. Where writers are
concurrent the question is real, and it has two answers rather than one,
because the obvious one does not scale to a traversal.

A slot the operation genuinely depends on can be entered into its read set as a
guard, checked at the commit point --- the model's opt-in, per-slot read
validation (\cref{sec:model}). But a guard is not a passive note in a side
table. It is a record like any other, and it installs in the slot exactly as a
store does, so a peer that writes that slot fails its own install and retries.
That is what makes it a guarantee rather than a hope, and it is why it costs
what a write costs. Guarding \emph{every} pointer a traversal crossed would
therefore be close to the literal reading of the question: the walk would leave
a proxy on every slot it stepped through, and every peer that touched any of
them would abort. An operation guards the few dependencies it has, for the same
reason it does not write words it has no business writing.

The traversal itself takes the other route, and it is a coarse one. It stays
ordinary uninstrumented reads, cheap and unguarded, and the \pseudotxn commits
optimistically; where peers keep invalidating it, it escalates. Past a retry
budget scaled by its loads as well as its records, a handle takes a per-domain
FIFO-fair lane, every writer in that domain funnels through the lane behind it,
and the whole operation --- traversal and mutation together --- re-runs inside
that turn. The guarantee the traversal gets is the turn, not a proxy on each of
its steps; and the budget counts the loads precisely because the traversal, not
the write set, is most of what the lane would serialize. Both routes belong to
the companion paper.

\emph{Declaring a write set disjoint is a promise the facility keeps you to.}
At the other end, an embedder whose slots are distinct \emph{by construction}
can skip the find altogether by declaring the handle disjoint --- which is to
say, by giving up read-your-own-writes, since the declaration asserts there is
nothing to reconcile and the facility then blind-appends. It is the right
trade for a write set with no self-overlap and no walk over its own pending
state, and the wrong one everywhere else. The declaration is a contract on slot
\emph{addresses}, not on the
logical items an edit names: two edits the caller believes distinct can still
land on one slot --- the adjacent deletes above both re-point \code{P->next} ---
so it is safe only where the write set is disjoint by construction, not merely
expected to be. Break the promise and the fused reconciliation is skipped, the
earlier same-slot edit is lost, and the commit reports success ---
which is why a debug build
(\code{-DURCU\_TXN\_SW\_DEBUG\_DISJOINT}) traps a violated declaration, and why
the declaration should be made only where disjointness is structural rather than
merely observed.

An STM offers no such dial, and the difference is the bargain of this paper in
miniature. Its read barrier sits on \emph{every} read, so a search inside a
transaction is reconciled whether or not the caller gave it a thought --- and
every read pays, whether or not it needed to. Here reconciliation follows the
accessor. A writer that must walk its own pending state loads transactionally
and pays the find; one that need not, or that knows its slots are disjoint, pays
nothing; and a \emph{reader} is not in the scheme at all, since its traversal is
uninstrumented plain \rcu, indistinguishable from a read outside any
\pseudotxn{} --- which is why the reader carries none of the per-element state
and none of the read-path work \cref{sec:costcomparison} charges the
alternatives. The writer chooses per load what the STM applies wholesale, and
the reader is what the choice buys.
% -----------------------------------------------------------------------
% NOTE (drafting; retained deliberately -- see README framing rule).
% This section was rewritten when the sw facility gained read-your-own-writes
% (urcu-txn-dev a797a2eb, 89e8d822, f6e48016). BEFORE that,
% the sw facility had no transactional loads: composition was slot-disjoint ONLY,
% and adjacent deletes silently corrupted -- which was true at the previous pin
% (ec621268) and is the behaviour earlier drafts described.
%
% Do NOT reinstate the old "the concurrent facility's RYW is what it is FOR / the
% reason it has transactional loads at all" framing anywhere (it was in this
% NOTE and in the Sec. 10 P2 hook). It is now FALSE: BOTH facilities have RYW, so
% RYW is part of the SHARED programming model (P3), not P2's differentiator.
% What still distinguishes the concurrent facility is the machinery exclusion
% removes -- the install CAS, conflict detection, abort -- for MULTIPLE writers
% (helping/stealing were built and REMOVED; see the retired-variants note, and
% do not reintroduce them as current MW machinery). RYW is orthogonal to that
% and now lives on both sides.
% -----------------------------------------------------------------------

\subsection{Composing with reclamation}
\label{sec:reclaim}

The wrappers inherit the discipline of \cref{sec:background} without amendment,
because the commit is the publish step.

A node's payload is initialized before the commit that links it, and the commit
is the release edge, so the payload is published safely. Nodes that a committed
\pseudotxn unlinked were reader-visible, so they are freed by
\code{call\_rcu()} after a grace period. Nothing an aborted attempt allocated
was ever published, so it is freed immediately --- but only what \emph{that
attempt} allocated, never anything already reachable.

The proxies themselves are subject to the same rule, and it is why settle exists
(\cref{sec:lifecycle}): a proxy is reader-reachable while installed, so it may
be reclaimed only once settle has rewritten every slot that named it, and then
only after a grace period.
% -----------------------------------------------------------------------
% NOTE (drafting; retained deliberately -- see README framing rule).
% RESOLVED -- do not add the "no re-link before a grace period" trap here.
% It is MULTI-WRITER ONLY and belongs in the companion paper. Checked: rcu-
% txn-list.h (the concurrent list) warns that a deleted node must not be re-
% linked before a grace period, because "the deleting transaction may still
% hold a parked proxy in that node's next slot, awaiting its settle. A plain
% store into a proxied slot is then overwritten when that settle converts
% the proxy, and the write is silently lost -- a permanently incoherent
% edge." That hazard needs a PEER transaction holding an unsettled proxy.
% Under exclusion the sole writer's transactions are sequential and settle
% completes (lifecycle step 4) before the next begins, so no proxy of a
% previous transaction survives into the next one. Confirmingly, rcu-txn-sw-
% list.h carries no such warning: it mentions grace periods only for
% reclaim. Stating it here would import a multi-writer hazard into a single-
% writer paper and imply a restriction the sw facility does not have.
% -----------------------------------------------------------------------

% ===========================================================================
\section{Measured Cost}
\label{sec:measured}

% ---------------------------------------------------------------------------
% NOTE (drafting; retained deliberately -- see README framing rule).
% SCOPE, revised 2026-07-26. This section measures the READER AXIS only: zero
% writers, and ONE writer. There is deliberately NO writer-scaling panel -- the
% facility is single-writer by contract (sec:exclusion), so N writers would
% serialise on a mutex and measure the mutex, not the flip-latch. Concurrent
% writers are the MW engine's story and the evaluation paper's.
%
% Restricting to the reader axis is what admits the cross-scheme comparison in
% tab:readclass: on this axis the comparison is sound in the direction that
% matters (they charge readers MORE and deliver MORE), and the MV-RLU figure is
% nearly independent of its clock variant, since a reader reads the clock once
% per section and mvrlu_deref only compares against that cached value.
%
% TRAVERSAL ORDER IS LOAD-BEARING, do not "simplify" the two txn arms into one.
% Order is worth ~15% here, an order of magnitude more than the per-dereference
% cost being measured. An earlier version of this table compared the transacted
% list's forward+backward walk against rculist's forward+forward and showed the
% facility ~8% AHEAD of plain RCU, which was an artifact of the access pattern.
%
% Every figure is the mean of two 3 s runs from scripts/p1_readers.csv in the
% benchmark tree, generated by scripts/run_p1_readers.sh, engine b3e23f9f. The
% superseded scripts/p1_sw_list.csv was taken on engine 5fc3cadf with the
% mismatched traversal and must not be cited. Profile figures are from a cycle
% profile of a single-writer, no-reader run. Cite no percentage not in those.
% ---------------------------------------------------------------------------

\Cref{sec:costcomparison} argues on mechanism that this facility charges the
reader less than the alternatives, and \cref{sec:whyweakening} argues that the
charge it does levy falls on the writer. This section measures both claims on the
bidirectional list of \cref{sec:swlist}, along the axis the facility is built
for: readers scaling to 192 cores, against zero writers and against one.

\paragraph{Setup.} Two AMD EPYC 9654 processors, 96 cores per socket, 24 NUMA
nodes. The harness pins worker $i$ to physical core $i$ and every sweep stops at
192, so no measurement below places two threads on one core. Userspace \rcu's
QSBR flavor, \code{call\_rcu} with per-CPU worker threads, and jemalloc with
per-CPU arenas. The list holds 10\,000 nodes with a 200-node churn set. Each
figure is the mean of two three-second runs of the \code{bench\_list\_scale}
harness. The harness checks reader-visible coherence on every traversal it
performs, in both directions; it reported no violation in any run reported here.

\paragraph{Traversal order, and why two transacted arms.} The transacted list
walks forward and then \emph{backward} --- the traversal \cref{sec:onestore}
showed \code{rculist} cannot offer under mutation at all --- while
\code{rculist} can only walk forward twice. Both visit the same number of nodes,
but not in the same order, and order is not free: a backward second pass
re-touches the tail first, the hottest lines it just loaded, while a second
forward pass restarts at the head, the coldest. On this machine that difference
is worth about 15\%, which is an order of magnitude more than the
per-dereference cost this section exists to measure. Comparing across it would
measure the access pattern.

So the transacted list is measured twice. Against \code{rculist} it walks
forward twice, matching it exactly, and the ratio is then the resolve alone.
Against the schemes that \emph{do} offer a coherent backward walk it walks
forward and backward, which is the operation they share.

\paragraph{Reads pay nothing for the \pseudotxn layer.} \Cref{tab:readcost}
gives the read ceiling with no writer running, in millions of node visits per
second, both engines walking forward twice.

% NOT a float, deliberately. As a [tb] table this was placed at the top of the
% page the section starts on -- above the section heading, reading as a table
% belonging to sec:reclaim -- and \FloatBarrier did not hold it there. It is
% small, and it belongs beside the sentence that introduces it, so it is
% anchored. \captionof (caption, via subcaption) numbers it as a table.
\par\medskip\noindent\begin{minipage}{\linewidth}
\captionof{table}{Read ceiling, no writer running: millions of node visits per
second, \emph{both engines walking forward twice}. The two elements are
byte-identical --- 24 bytes, two pointers and an \code{int} --- so there is no
footprint difference and no access-pattern difference between the rows, and the
ratio is the resolve alone.}
\label{tab:readcost}
\centering
\begin{tabular}{@{}lrrrrr@{}}
\toprule
Readers & 1 & 8 & 32 & 96 & 192 \\
\midrule
\code{txn\_sw\_list} & 778 & 6\,250 & 24\,991 & 71\,774 & 146\,191 \\
\code{rculist}       & 761 & 6\,131 & 24\,518 & 71\,369 & 142\,367 \\
\midrule
Ratio                & 1.022 & 1.019 & 1.019 & 1.006 & 1.027 \\
\bottomrule
\end{tabular}
\end{minipage}\par\medskip

The two are indistinguishable: every ratio lies within 3\% of unity, and
run-to-run drift on this machine is itself about 0.8\%. Both scale essentially
linearly across the 192 cores --- the transacted list gains a factor of 188 from
one reader to 192, or 97.8\% of ideal, and \code{rculist} gains 187, or 97.4\%.

That the two agree is the point, and the construction is what makes it a fair
statement. The rows are not merely close; they are running the same traversal
over elements of the same size. \Cref{tab:comparison}'s zero in the loads column
predicts exactly this: the marker rides in the word the reader had to load
anyway, so the only residue is a predicted branch. We do not claim the small
positive margin --- the facility does strictly more work than plain \rcu, so a
real advantage would need a mechanism we cannot name, and the margin wanders
without pattern. The claim is parity.

\paragraph{What the stronger guarantees cost the reader.}
\Cref{tab:readclass} places the same traversal against the schemes that also
deliver a coherent bidirectional walk. All of them give readers strictly more
than this facility does --- a coherent snapshot, or per-element existence, where
a \pseudotxn gives only monotone resolution (\cref{prop:monotone}) --- so the
gaps are what that extra guarantee costs on the read side, which is the
comparison \cref{sec:whyweakening} promises and never quantifies.

\par\medskip\noindent\begin{minipage}{\linewidth}
\captionof{table}{Read ceiling, no writer, all engines walking forward then
backward: millions of node visits per second, and the ratio to the transacted
list. Existence is measured in its best per-element layout, carrying only its
tagged word (a 32-byte element) with the writer-only state segregated; charging
it for the reference implementation's 136-byte element would measure a struct
layout rather than the design. MV-RLU runs its global-clock configuration, which
needs no calibrated per-machine constant.}
\label{tab:readclass}
\centering
\begin{tabular}{@{}lrrrr@{}}
\toprule
Readers & 1 & 32 & 192 & Ratio \\
\midrule
\code{txn\_sw\_list}    & 788 & 25\,339 & 148\,190 & 1.00 \\
Existence               & 691 & 22\,215 & 123\,035 & 1.14--1.21 \\
RLU                     & 622 & 19\,943 &  99\,017 & 1.27--1.50 \\
MV-RLU                  & 293 &  9\,165 &  46\,699 & 2.69--3.23 \\
\bottomrule
\end{tabular}
\end{minipage}\par\medskip

Existence is the closest, and the residue is precisely what
\cref{tab:comparison} attributes to it: one acquire load per dereference, and a
32-byte element against 24, so fewer elements share a cache line. That the
margin is only 14--21\% is the honest form of the claim, and a stronger one than
a larger number would be --- it says the \pseudotxn wins because it attaches
\emph{no} per-element state, not because an implementation chose a fat one.

\paragraph{The writer pays about twice.} On the same list, with a single writer
and readers running concurrently:

% Anchored, for the same reason as tab:readcost above.
\par\medskip\noindent\begin{minipage}{\linewidth}
\captionof{table}{Single-writer update throughput with concurrent readers,
millions of operations per second, each engine under an uncontended writer
mutex. Both sides reclaim through \code{call\_rcu()} and neither changes \rcu
flavor, so the ratio isolates the \pseudotxn layer.}
\label{tab:writecost}
\centering
\begin{tabular}{@{}lrrrrr@{}}
\toprule
Readers & 1 & 8 & 32 & 96 & 191 \\
\midrule
\code{rculist}       & 14.25 & 4.02 & 3.93 & 3.21 & 2.92 \\
\code{txn\_sw\_list} &  7.12 & 3.05 & 2.71 & 1.71 & 1.27 \\
\midrule
Ratio                & 2.00 & 1.31 & 1.45 & 1.88 & 2.30 \\
\bottomrule
\end{tabular}
\end{minipage}\par\medskip

The tax is about twice, and it is not a constant: it falls to 1.31 at eight
readers before climbing back to 2.30 at 191. The shape is informative. Plain
\rcu's writer collapses from 14.25 to 4.02 between one reader and eight, while
the transacted writer only falls from 7.12 to 3.05 --- grace periods lengthen
for both, and the cheaper writer has more to lose. They converge under reader
load and diverge again as the grace period stretches further.

With no reader at all the ratio is wider, 2.84 under the same mutex. We report
the figures with readers because that is the regime the facility is for: a
structure with no readers has no use for reader-visible coherence.

\paragraph{Where the writer's cost goes.} A writer tax of two to three times
invites the conclusion that the flip is expensive. It is worth disposing of
that, because the profile says the opposite. In a cycle profile of a
single-writer run with no reader --- the widest ratio, and so the least
favorable place to look --- the commit
itself, \code{urcu\_txn\_sw\_commit\_flavor}, accounts for 3.4\% of cycles. Ahead
of it sit the allocator, at 10.7\% in jemalloc's \code{malloc}, and the per-CPU
\code{call\_rcu} plumbing: 7.7\% in \code{sched\_getcpu} and 8.4\% across the two
lookups that find the current CPU's queue. Those are \rcu's shared
infrastructure, and plain \rcu pays them too --- just fewer times.

The mechanism sits in the reclaim contract rather than in the commit. A list
operation moves two edges, so it parks two proxies, and a parked group block must
outlive a grace period before it can be freed, because a reader that loaded a
proxy may still be resolving through it (\cref{sec:reclaim}). Every update
therefore defers a second object --- the descriptor --- in addition to the node it
unlinked, where plain \rcu defers only the node. A commit that moves a
\emph{single} edge needs no proxy and frees in place. So the price of a coherent
backward chain, in this implementation, is largely one extra
grace-period-deferred free per update.

Two of those costs have known remedies, both in \rcu's own \code{call\_rcu}
machinery rather than in this facility: an \code{rseq}-supplied CPU id would
remove the
\code{sched\_getcpu} call, and an \code{rseq}-based enqueue would shorten the
rest of the \code{call\_rcu} path. Both would help plain \rcu equally, so neither
narrows the ratio so much as lowers both sides. The extra deferred free is
intrinsic to publishing two edges atomically, with one exception the facility
already admits: an embedder that knows no reader can hold a proxy may pass a
synchronous reclaim shim and free in place.

\paragraph{What this does not measure.} Three absences are deliberate.

There is no writer-scaling panel. The facility is single-writer by contract
(\cref{sec:exclusion}), so running $N$ writers against it would serialize them
on a lock and measure that lock. Concurrent writers are a real question and this
is not the design that answers them; they belong to the multi-writer engine and
to the evaluation paper. The reader axis is where a read-mostly facility earns
or loses its keep, and it is the axis measured here.

The cross-scheme figures in \cref{tab:readclass} are read-side only, and are not
a throughput evaluation of RLU, MV-RLU or existence. Their write sides differ in
ways this workload does not isolate --- MV-RLU reclaims through a single
background thread, existence and this facility defer a group object per update
--- and the guarantees differ, so a write-side ranking across these rows would
compare unlike things. What the reader axis supports is the narrow claim made
above: what the stronger reader guarantee costs the reader.

And this workload cannot separate read-side synchronization costs, which is worth
stating plainly because the same harness reports them. A \emph{reader-biased}
reader-writer lock and a seqlock traversing the same list land within a few
percent of plain \rcu, and at 192 readers the seqlock is the fastest of the four.
That is not a finding about locks --- least of all a flattering one, since a
reader-preferring lock is exactly the variant a read-only workload suits. It is
what happens when a 10\,000-node traversal amortizes one acquisition to nothing,
and it is why the read-side comparison in
\cref{tab:readcost} is against \code{rculist} alone. Separating those costs needs
a cache-resident list or a concurrent writer, and belongs to the evaluation
paper along with everything else this section declines to claim.

% ===========================================================================
\section{Limitations}
\label{sec:limitations}

The facility's limits divide cleanly in two, and the division is the useful part.
Some are \emph{properties of the model} --- they cannot be violated, only
disliked. The rest are \emph{contracts on the embedder} --- they can be
violated, and violating them is expensive.

\subsection{What holds unconditionally}
\label{sec:unconditional}

One thing is worth stating positively first, because it is the guarantee that is
usually hardest to get in this design space and it costs nothing here.

The facility is indifferent to \emph{slot-value} A-B-A. Resolution goes through
the record's own status rather than through the slot's plain value, so the
identity that matters is the descriptor's, not the value's. Any slot value may
recur --- a next-pointer cycling $B \rightarrow X \rightarrow B$ with $B$ a live
successor that is never freed, a counter revisiting a number, an embedder's own
reuse --- and the recurrence is inert. \rcu prevents \emph{address} reuse, but
not \emph{value} recurrence; this facility does not need it to.
% -----------------------------------------------------------------------
% NOTE (drafting; retained deliberately -- see README framing rule).
% Be precise about the scope of "unconditional" here, because it is easy to
% overstate and the overstatement is load-bearing. The facility's own words
% are "memory safety, which holds unconditionally here" -- and the sentence
% is about SLOT VALUES. It is NOT a claim that memory safety survives a
% violated contract. It cannot be: Sec. 7.2's tag contract, violated, makes
% a reader follow a non-proxy value as a record address. Do not write
% "memory safety is unconditional, atomicity is conditional" as a blanket --
% it is wrong in the direction that flatters us, which is the worst
% direction.
% -----------------------------------------------------------------------

\subsection{The contracts, and what violating them costs}
\label{sec:contracts}

The facility requires four things of its embedder. None is checked in a default
build.

\begin{description}
\item[Writer exclusion.] The mechanism has no install CAS, no conflict
  detection and no abort (\cref{sec:exclusion}). Two writers racing one slot
  corrupt it.
\item[Free tag bits.] Every value a transacted slot may legitimately hold must
  satisfy \code{(value \& tag) != tag} (\cref{sec:tagcontract}).
\item[Correct composition, and no raw access.] Edits that touch the same slot
  must be composed through the read-your-own-writes path, which fuses them
  (\cref{sec:swcomposition}); the blind fast path and any handle declared
  disjoint instead require pairwise-distinct slots. Every write to a transacted
  slot must go through the facility and every read through its accessor.
\item[Deferred descriptor reclaim.] A descriptor reachable through a slot is
  reclaimed only after a grace period, because a reader may dereference it.
\end{description}

What a violation costs is worth stating exactly, because it is more than
atomicity. A value that accidentally carries the tag pattern is
indistinguishable from a proxy, and a reader will untag it and dereference it as
a record: that is a wild pointer, not a torn read. Racing writers corrupt slots
whose values readers then follow. \textbf{These contracts underwrite memory
safety, not merely atomicity} --- and that is the honest shape of the safety
argument.
% -----------------------------------------------------------------------
% NOTE (drafting; retained deliberately -- see README framing rule).
% This is the series' central honesty point and Sec. 1 should not oversell
% around it. The July review of this facility concluded that every major
% finding sat at the CONTRACT layer rather than the mechanism layer -- doc
% text with API force steering embedders into corruption. That is a real
% result about this class of design and worth saying: the mechanism is the
% easy part.
% -----------------------------------------------------------------------

\subsection{Contracts the facility cannot check for free}
\label{sec:checkability}

The contracts of \cref{sec:contracts} are not equally checkable, and the facility
is candid about which is which rather than uniformly optimistic.

Exclusion has an opt-in runtime validator
(\code{-DURCU\_TXN\_SW\_EXCL\_VALIDATE}) that aborts on violation and costs
nothing when off --- the handle does not even carry the owner field. It claims
the handle (catching a \pseudotxn handed between threads) and the slot (where
two racing writers actually collide).

The tag contract has no such check, and cannot cheaply have one: verifying it
would mean inspecting every value ever stored in a transacted slot, which is the
per-write cost the design exists to avoid. It is discharged by the embedder's
encoding, by construction or not at all.

A disjoint declaration (\cref{sec:swcomposition}), if false, is the quietest
violation of all: the facility skips the reconciliation the declaration promised
was unnecessary, drops the earlier of two edits that land on one slot, and
commits success. It is the one contract with a dedicated opt-in trap
(\code{-DURCU\_TXN\_SW\_DEBUG\_DISJOINT}), though even that catches only the
recording side; a default build discharges the declaration by construction or
corrupts quietly.
% -----------------------------------------------------------------------
% NOTE (drafting; retained deliberately -- see README framing rule).
% Frame the opt-in validator as the design position it is: the contract is
% unenforceable at zero cost, so the facility offers a check that is free when
% off rather than pretending the contract is safe or taxing everyone to
% police it. That is the same trade as the rest of the paper -- pay nothing
% by default, and be explicit about what that costs.
% -----------------------------------------------------------------------

\subsection{Properties of the model, not defects}
\label{sec:modellimits}

The remaining limits cannot be violated because they are not promises. They are
what the model \emph{is}.

\textbf{No reader snapshot.} A reader is not a transaction of any kind
(\cref{sec:nonguarantees}); a traversal spanning slots may straddle a commit. An
embedder that needs a coherent multi-slot read layers its own versioning on top.
Two standard families do this without disturbing the base facility: a
\emph{seqcount}, where the reader re-samples a counter around its read and
retries if a writer intervened, and a \emph{validating transaction}, where the
reader's loads fold into a commit that is accepted only if every slot it touched
still holds what it returned --- the loads ``act as writes''. Their progress
guarantees and cost are the subject of later work in this series. That the base
model supplies neither is the deliberate weakening of \cref{sec:whyweakening},
not a shortfall.
% -----------------------------------------------------------------------
% NOTE (drafting; retained deliberately -- see README framing rule).
% DELIBERATELY UNDERSTATED. This bullet names only the two STANDARD families
% (plain seqcount; validating transaction) and stops. It does NOT describe how
% far the read side can be kept wait-free, and must not.
%
% Reason (Mathieu's call): the single-bit grace-period-gated seqcount latch with
% a copy-on-write overflow escape -- which makes a torn-free multi-word read
% ALSO wait-free (<= 1 retry, constant in writer count) by admitting <= 1
% in-place update per GP per node and recompacting the rest -- is suspected
% NOVEL and is RESERVED for its own paper (a fuller, enabling disclosure is
% stronger prior art than a buried hint, and self-disclosure never blocks the
% later paper). Mechanism lives in efficios-trie-benchmark design/rcu-txn-blob.md
% ("Single-bit latch: wait-free reads for real"); the urcu-txn-dev commits that
% carry it are unpushed and volatile (hashes change on amend) -- do not cite them.
%
% So: do NOT reinstate the wait-free single-bit-latch mechanism, its retry bound,
% the COW escape hatch, or any "wait-free XOR torn-free / COW is the only
% wait-free corner" framing here. An earlier draft did (following the stale
% "seqcount reads are not wait-free" note) and it was both wrong AND a premature
% disclosure of the reserved result. Keep P1 neutral on read-side progress class.
% -----------------------------------------------------------------------

\textbf{Value-based, not stability-based.} A record's old is the value the slot
must hold \emph{at} the commit point, not one it must have held throughout. An
operation whose correctness requires \emph{detecting} an away-and-back change
must carry its own version in the word. Under exclusion the requirement does not
arise: the one thing that could invalidate a recorded old --- an earlier edit
the writer staged in the same \pseudotxn{} --- is exactly what read-your-own-writes
reconciles (\cref{sec:swcomposition}), and no other writer can change the slot
at all. The distinction earns its keep only where writers are concurrent and a
peer can move a slot away and back between a read and the commit; under
exclusion it is inert.
% -----------------------------------------------------------------------
% NOTE (drafting; retained deliberately -- see README framing rule).
% UPDATED when the sw facility gained read-your-own-writes. The
% earlier draft unified "value-based not stability-based" with "composition is
% slot-disjoint only" -- true when the facility could NOT read its own pending
% edits, false now that it can. Under exclusion the only thing that could
% invalidate a recorded old is your own earlier same-bracket edit, and RYW now
% reconciles exactly that, so the value/stability distinction is inert here
% rather than surfacing through composition. It earns its keep only in the
% concurrent facility, where a PEER can change a slot between read and commit --
% which is where the guard/validate API lives. Do not re-couple it to
% composition here.
% -----------------------------------------------------------------------

\textbf{Monotonicity is a property of the traversal, not only of the facility.}
\Cref{prop:monotone} holds for a dependency-chained reader. One that hops
without a chain --- re-reading a cached pointer, or comparing two
independently-reached slots --- must supply its own acquire
(\cref{sec:monotone}). The selector's own monotonicity is unconditional; what a
reader's sequence inherits from it is not.

% ===========================================================================
\section{Variations and Alternatives}
\label{sec:variations}

% ---------------------------------------------------------------------------
% THIS SECTION IS THE DEFENSIVE PAYLOAD. A patent claim can cover an obvious
% variant of what is disclosed, so the design space NOT taken must be disclosed
% and dated too. A normal paper cuts this; here it is load-bearing.
%
% RULE: every variant must be ENABLING -- described well enough that a person
% skilled in the art could build it. A variant named but not explained
% discloses nothing.
%
% Do NOT announce the defensive motive in the text. The section stands on its
% academic merit (readers genuinely learn the trade-offs), and the disclosure
% works regardless of stated intent. Announcing it invites hostility and buys
% nothing. If Mathieu wants the motive stated, that is his call -- but the
% default is silence.
%
% The MULTI-WRITER retired variants -- helping, stealing, the per-record install
% latch, invisible writes, the k=1 Bloom filter, the flat budget, the dense-array
% RYW overlay -- belong to P2. They were built, measured and removed, and that
% falsification trail is P2's strongest asset. Do not spend it here.
% ---------------------------------------------------------------------------

The mechanism of \cref{sec:swmechanism} is one point in a design space. This
section maps the space around it: the choices available at each decision, what
each would cost, and why this facility chose as it did. Several of these
alternatives are perfectly workable; the reasons for rejecting them are
quantitative or contextual, not fundamental.

\subsection{The selector}
\label{sec:varselector}

\paragraph{Boolean versus generation counter.} The selector here is a boolean
written once, $0 \to 1$. A generation counter is the obvious alternative: each
proxy records the generation at which its new value becomes current, and
resolution compares the group's counter against it. That buys reusable groups
--- a group could serve many successive transactions without reallocation ---
and it is how clock-based schemes such as RLU decide visibility. It costs the
write-once property, and with it the unconditional monotonicity of
\cref{prop:monotone}: a counter that advances repeatedly makes ``never
new-then-old'' a claim about a range rather than a single transition, and a
reader must then compare rather than test. We chose the boolean because
\cref{prop:monotone} is what lets embedders trade a snapshot for a fixed read
order (\cref{sec:publishorder}), and a boolean makes it free.

\paragraph{One word per group versus packed.} Each group here owns a full word.
Selectors could instead be packed --- many groups' bits in one word, with each
proxy carrying a bit index --- which shrinks the group to nothing but makes the
selector word a contention point shared across unrelated transactions, and makes
each flip a read-modify-write rather than a store. Under exclusion the
read-modify-write is safe; under concurrent writers it would not be.

\paragraph{Write-once versus reusable.} A group here is used once and reclaimed.
It could be reset and reused, which trades the reclamation cost for a
quiescence problem: the reset must not overtake a reader still resolving a proxy
of the previous use, so the reuse needs exactly the grace period the
reclamation already pays.

\subsection{The tag}
\label{sec:vartag}

\paragraph{Width, and the alignment it implies.} A one-bit tag needs only
two-byte alignment of the record, which any pointer has. The four-bit tag this
facility permits requires sixteen-byte alignment, which must be arranged
deliberately (\cref{sec:alignment}). Wider tags are available at proportionally
stricter alignment: an eight-bit tag would need 256-byte-aligned records, which
is affordable for a small pool of proxies and wasteful for many. The width an
embedder needs is set by its own encoding, not by the facility.

\paragraph{Low bits versus high bits.} This facility takes tag bits from the low
end, where alignment guarantees them. The high end is the alternative: on
\emph{most} current 64-bit architectures the top bits of a virtual address are
unused --- though not all, s390x being an exception --- and a tag can live there
instead, leaving the low bits entirely to the embedder.
The first cost is portability --- the number of free high bits is
architecture-dependent and can change --- so a high-bit scheme needs a runtime
probe of the address space to establish its budget, and a fallback for
architectures that offer fewer bits than it wants.

The second is that the reader's test is not equally cheap at both ends, and on
x86-64 the asymmetry falls the way the tag's width decides. A low mask is
an immediate operand: one \code{test} against an immediate, whatever the reader
already holds. A mask of several \emph{high} bits does not fit that immediate, so
it must be materialized into a register first --- an extra instruction and a
register the reader did not otherwise need, on the fast path, per dereference. A
\emph{single} high bit escapes this, because a bit-test instruction takes the bit
index as a small immediate rather than a mask. So the reserved nibble of
\cref{sec:alignment} is the case that pays for living high, and a lone tag bit is
close to free at either end. Like the budget itself, this is
architecture-dependent: an instruction set with richer logical immediates prices the two ends
differently.

Both of those concern the \emph{test}. Recovering the proxy is a separate
operation, and there the asymmetry does not soften. This facility recovers it by
\emph{subtracting} the tag rather than masking it off, which is exact because the
two ends of the encoding meet: a value is untagged only once the test above has
proven every tag bit set in it, and a proxy's own address has every tag bit
clear, so the tag is precisely the difference. With a compile-time-constant tag
the subtraction then folds into the addressing-mode displacement of the load that
follows, and the head of the resolve's dependency chain issues straight off the
raw tagged value --- where a mask cannot fold, and sits as an ALU operation
between the slot's load and the first load that depends on it.

The subtraction is the better instruction even where it does not fold, because it
leaves the proxy at a constant offset from the value the reader loaded, which
address generation can follow and a mask's bitwise rewrite does not. A high tag
gets none of this: too large for a displacement, it needs an instruction of its
own ahead of that load even as a single bit. All of which lands on the taken
path, and \cref{sec:costcomparison} counts that path as bounded rather than
steady state --- so this decides nothing on its own. It is offered because the
same reasoning that puts the tag in the low bits pays a second time there.

C++ standardization has come down on the same side. The pointer-tagging proposal
noted in \cref{sec:alignment} exposes the low bits only, deriving how many are
available from the pointee's alignment, after SG1 declined to standardize access
to the high bits as dangerous and non-portable and as a constraint on the
security features that may yet claim them~\citep{wg21p3125}. That is a narrower
judgement than ours --- it is about what a language should promise, not about
what an embedder may do on a platform it knows --- but it lands on the same
asymmetry: the low bits are the ones an implementation can be held to.

The high route is not hypothetical either, and what it costs in practice is worth
knowing. ZGC colors every object reference with four metadata bits placed above
the address bits, and interprets them in a load barrier on every reference load
\citep{yang2022zgc}. That is a working high-bit tag in a production collector ---
and it is also the read-side bargain \cref{sec:whyweakening} declines, since the
barrier runs whether or not any relocation is in flight. The comparison cuts both
ways: the high bits are usable, and using them for a marker a reader must
interpret puts work back on the reader, which is the cost this facility exists to
avoid wherever the marker lives. A facility could still offer both and let the
embedder choose.

\paragraph{Tagging versus a sentinel range.} Tagging is not the only way to make
a proxy recognizable. An embedder could reserve a range of values --- a page of
address space, say, or an unused numeric range --- and treat any value in it as
a proxy handle. This costs no bits of the word, which is the one thing tagging
does cost, and it is a real option for a slot whose bits are all spoken for. It
costs a heavier check on the read path --- typically two inequality comparisons
to bracket the value against the reserved range, where the tag is a single
mask-and-compare --- plus an allocator constrained to that range, and a carve-out
from the data domain that the embedder must guarantee is never a legitimate
value.

Anchoring the range at zero recovers most of that read-path cost. If a handle is
an \emph{index} into the reserved pool rather than an address within it, the
check is one unsigned comparison against the pool's size --- as cheap as the
tag's mask-and-compare, and the carve-out becomes the low addresses no allocator
returns anyway. What anchoring does not change is the pool. The allocator behind
it must be a real one, not a fixed table indexed by the handle, because a proxy
must outlive a grace period: a reader that loaded a handle before the commit
settled may still be resolving it, and a table sized to a \pseudotxn's width
would recycle that index under such a reader --- silently, since the recycled
index resolves to a well-formed proxy belonging to a different transaction. The
bound is therefore not the widest transaction but the number of proxies live
within a grace period, which is a rate rather than a shape: a lone writer
committing steadily leaves every settled descriptor lingering until the readers
that might hold a handle have gone quiet, so the pool must cover the update rate
multiplied by the grace-period duration. Exclusion buys nothing here --- the
lifetime being waited out is a reader's, not a writer's. A range allocator over
a pool sized that way keeps every proxy independently reclaimable, which is what
lets many transactions be in flight, and lingering, at once.

\subsection{Proxy placement and allocation}
\label{sec:varproxy}

\paragraph{One block versus many.} A proxy is a record in its installed form
(\cref{sec:fliplatch}): the tagged object a slot points at while the \pseudotxn
is live. Proxies and their group are typically allocated as a single backing
block with one \code{rcu\_head}, so the whole set costs one allocation and one
deferred free. Individually allocated proxies are
the alternative: they cost an allocation and a reclamation each, and they permit
a write set whose width is not known in advance to grow without moving what it
has already built --- which is the problem \cref{sec:tagnotptr} solves the other
way, by making a move safe.

\paragraph{Referenced versus embedded.} A slot here points at a proxy. A proxy
could instead be embedded in the target object, so the slot points at the object
as usual and the object's header says whether an indirection is in flight. That
is not a hypothetical design. It is the forwarding pointer of
\citet{brooks1984forwarding}: a word in every object's header naming either the
object itself or its relocated copy, which a reader follows
\emph{unconditionally} rather than testing --- and it shipped that way in
Shenandoah's concurrent compacting collector for OpenJDK
\citep{flood2016shenandoah}.

The trade is the one \cref{sec:whyweakening} declines: a header word on every
object for as long as it lives, and an indirection on the read path whether or
not any mutation is in flight. That is a good trade in a collector, where any
object may have to become relocatable at any moment and the barrier must
therefore be unconditional anyway. A \pseudotxn needs the indirection only while
a commit is in flight on a particular slot, which is why it puts the indirection
in the slot instead of the object and takes it away at settle.

\subsection{Settle}
\label{sec:varsettle}

\paragraph{Eager, lazy, or never.} This facility settles eagerly: after the flip,
every slot is rewritten from the tagged proxy back to the direct target. Settle
could be deferred, batched across transactions, or omitted entirely.

Deferral is the interesting middle, and what limits it is not that two
transactions may want the same slot. Under exclusion that costs nothing to
resolve: a writer installing into a slot its own earlier transaction left
unsettled reads through the stale proxy for its old value, exactly as a reader
would, overwrites the slot with its own proxy, and drops the pending settle ---
a local fact, recorded by the only thread entitled to touch that slot at all.
What limits deferral is that batching moves the settle out of the region the
exclusion covers. Settle is a plain store only because the caller holds a lock
over the slot while it runs; accumulate settles across transactions and the
flush happens once that lock is gone, over slots that may be covered by
different locks. The batch must then re-acquire each of them --- paying back the
amortization it was collecting, and inheriting a lock-ordering problem the eager
form never has --- or make each settle conditional on the slot still naming its
own proxy, which turns the \sw path's plain store into an atomic and gives up
precisely what makes that path cheap. Deferral also keeps proxies installed for
longer, so more descriptors are reader-reachable at once, on the same
rate-times-window budget as \cref{sec:vartag}.

Omitting settle entirely is the endpoint, and it is \emph{correct}: the proxies
resolve correctly forever.
It is also unbounded: every read of that slot pays an indirection permanently,
the proxies can never be reclaimed, and the structure accumulates a growing
layer of resolved-but-still-present latches. Settle exists to make the cost
transient, which is the whole thesis of \cref{sec:whyweakening}; a facility that
never settled would have per-slot state after all, arrived at by a different
road.

\subsection{Reclamation}
\label{sec:varreclaim}

Proxies are reader-reachable while installed, so their reclamation needs a
scheme that knows when readers have finished. \rcu is one; the others are the
usual suspects and they are usable here. \code{synchronize\_rcu()} instead of
\code{call\_rcu()} makes the writer wait rather than defer, which is simpler and
serializes the writer against every reader in the system. Hazard pointers
\citep{michael2004hazard} would let a proxy be freed as soon as no reader holds
it, at the cost of a per-read store --- which is exactly the read-side charge
this design refuses (\cref{sec:whyweakening}), so it would be a strange choice
here while being a reasonable one elsewhere. Reference counting on the group
would serialize readers on a shared counter.
% -----------------------------------------------------------------------
% NOTE (drafting; retained deliberately -- see README framing rule).
% This subsection matters more than it looks for disclosure purposes: it
% covers the "same mechanism, different reclamation" variants. Keep the
% reasoning attached to each -- "hazard pointers would work but cost a per-
% read store" is enabling; "we use RCU" is not.
% -----------------------------------------------------------------------


% ===========================================================================
\section{Comparison with prior art}
\label{sec:costcomparison}

\Cref{sec:whyweakening} argued that every design adding multi-pointer atomicity
to \rcu must charge the reader something, and that this facility charges the
least. This section makes the comparison concrete. \Cref{tab:comparison}
tabulates the read-side cost each design pays per dereference; the discussion
that follows places existence structures, RLU and its multi-versioned successor,
and software transactional memory against this work in turn.

\begin{table}[tb]
\centering
\footnotesize
\setlength{\tabcolsep}{4pt}
\begin{tabular}{@{}l l c c l l@{}}
\toprule
 & \textbf{Per-element} & \textbf{Extra} & \textbf{Extra} & \textbf{Reader} & \textbf{Publish} \\
 & \textbf{state} & \textbf{loads} & \textbf{branches} & \textbf{guarantee} & \textbf{step} \\
\midrule
Existence~\citep{mckenney2016beyondissaquah}
  & \code{eh\_egi} field
  & 1
  & 1
  & per-element
  & group state \\
RLU~\citep{matveev2015rlu}
  & object header
  & 1
  & 1
  & snapshot
  & global clock \\
MV-RLU~\citep{kim2019mvrlu}
  & header $+$ versions
  & $\geq 2$
  & $\geq 3$
  & snapshot
  & global clock \\
\midrule
\textbf{This work}
  & \textbf{none}
  & \textbf{0}
  & 1
  & never new-then-old
  & selector \\
\bottomrule
\end{tabular}
\caption{Where each design pays for multi-pointer atomicity, on the read-side
fast path, per dereference. Each publishes with a single store to the shared word
named in the last column, so publish is one store for every design, not only
this one. \emph{Extra loads} counts additional memory
references over and above the load the reader must issue anyway; \emph{extra
branches} counts conditional branches. Every design tests something, this one
included --- the honest difference is in the loads column, not the branches
column. Existence's load is an \emph{acquire} of a per-element field; RLU's is a
plain load from its per-thread state. A \pseudotxn's is absent because its
marker rides in the value \code{v} that \code{rcu\_dereference} already
returned, so its test is arithmetic on a register the reader already holds ---
but the test is still a branch, and \textbf{0} in the loads column means zero
loads, not zero instructions. That \textbf{0} is the no-update figure: a reader
that meets a commit in flight on the slot resolves it through the proxy and
selector, the pair of loads this path avoids. All three single-branch entries are
predicted
branches, statically hinted in their respective sources, and each costs a
branch-predictor entry. The first column is the unconditional cost, and is where
the designs differ permanently: existence and RLU attach state to every element
for as long as the structure lives, whereas a \pseudotxn attaches nothing and
confines its state to the mutation. The existence row names what the mechanism
\emph{requires}, which is the one word the reader tests; a reader who consults
the reference implementation will find considerably more embedded beside it ---
group linkage, a reclamation head, callbacks and the writer's lock, 112 bytes in
all, of which only the 8-byte word is ever read on the traversal path. That
surplus is an implementation choice rather than a property of the design, and
it is separable: \cref{sec:costwalk} measures what separating it is worth.
MV-RLU is the one entry whose counts are not
constants: its dereference reads the object header unconditionally --- it has no
counterpart to RLU's pure-reader bypass --- and then may walk a version chain
whose length the reader does not control, so its figures are floors rather than
typical cases. These counts are \emph{measured}, not inferred from the sources:
hardware counters over a single-reader traversal, differenced against a longer
run so that startup cancels, give the marginal loads and branches per node
visit, and the figure quoted is each design's count minus plain \rcu's. Every
entry in the two count columns survives that check --- this work at exactly
zero extra loads and one extra branch, existence at one and one, RLU's
dereference at one --- and MV-RLU exceeds its floors by a wide margin, at ten
extra loads and fifteen extra branches.}
\label{tab:comparison}
\end{table}

% ---------------------------------------------------------------------------
% PROVENANCE OF EVERY CELL IN tab:comparison -- VERIFIED FROM SOURCE 2026-07-16.
% Kept as a LaTeX comment rather than a \NOTE because it quotes C: the bare _, &
% and ~ would need escaping inside a macro argument, and an escaping slip here
% breaks only the draft build, which is exactly when you need to read this.
% Do not "tidy" these cells from memory. What was actually read:
%
% existence_exists()  --  perfbook datastruct/existence/existence.h:116
%     uintptr_t egi = smp_load_acquire(&ehp->eh_egi);
%     if (likely(!egi)) return 1;
%     esp = (uintptr_t *)(egi & ~0x1);
%     return (egi & 0x1) == smp_load_acquire(esp);
% The FAST path is ALREADY one acquire load of a PER-ELEMENT field (eh_egi).
% When the element is in a group there is a SECOND acquire load, of the group's
% state word -- and struct existence_group is __attribute__((aligned(
% CACHE_LINE_SIZE))), so that is a line shared by every reader of every element
% in the group. Note also: existence uses pointer tagging in bit 0 of eh_egi
% ("Dmitry Vjukov's pointer-tagging trick", per its own comment). Say so.
% Tagging is not our invention and claiming it would be trivially refuted.
%
% STRUCT SIZE, verified with pahole 2026-07-26, because the "eh_egi field"
% cell understates what the reference implementation actually embeds:
%   struct existence_head = 112 bytes, 2 cachelines, 1 four-byte hole
%     eh_egi   0   8   <- the ONLY field on the read path
%     eh_list  8  16   group linkage, live only while a flip assembles
%     eh_add/eh_remove/eh_free 24 24   compile-time constants for an embedder
%     eh_gone 48   4   + 4-byte hole
%     eh_lock 56  40   spinlock_t, and perfbook/api.h:188 typedefs that to
%                      pthread_mutex_t -- so 40 bytes, straddling the 64-byte
%                      boundary. NOT our inflation; check api.h before doubting.
%     eh_rh   96  16   reclamation
% Embedded in a list element that also carries two pointers and a key, the
% reference layout puts the key at offset 128 -- cacheline 2 -- so a traversal
% touches TWO lines per element. That is worth ~2x (scripts/existence_layout.csv
% in the benchmark tree). The caption now says the surplus is separable and
% points at sec:costwalk; do NOT "correct" the first column to 112 bytes, which
% would charge the DESIGN for an implementation choice, and do not drop the
% clause either, which would leave a reader who greps existence.h thinking the
% row is wrong.
%
% RLU_DEREF_INTERNAL  --  third_party/rlu/rlu.h:221
%     if (rlu_likely(self->is_check_locks == 0)) { p_cur_obj = p_obj; }
%     else { if (p_obj && RLU_IS_UNLOCKED(p_obj)) ... else slow_path }
% and rlu.c:797 clears is_check_locks at reader-lock time when the thread has no
% pending writes of its own (n_pure_readers++). So a PURE READER never touches
% the object header: a thread-local load and a predicted branch, nothing more.
%
% DO NOT claim "RLU costs an extra load on every read". It is FALSE for pure
% readers and a reviewer who knows RLU will catch it on sight. RLU's
% unconditional cost is SPACE -- rlu_obj_header_t is prepended to every object
% (RLU_MOVE_PTR_BACK) -- not instructions. An earlier draft of this table made
% exactly that error; it was caught only by reading rlu.h.
%
% mvrlu_deref()  --  github.com/cosmoss-jigu/mv-rlu, lib/mvrlu.c, read 2026-07-16
%     p_act = get_act_obj(obj);                    /* reads the object header */
%     p_copy = vobj_to_obj_hdr(p_act)->p_copy;     /* second header load      */
%     if (unlikely(p_copy)) {
%         do {
%             chs = vobj_to_chs(p_copy);
%             wrt_clk = get_wrt_clk(chs);
%             if (lte_clock(wrt_clk, self->local_clk)) return p_copy;
%             if (unlikely(lte_clock(chs->cpy_hdr.wrt_clk_next, qp_clk2)))
%                     break;
%             p_copy = chs->obj_hdr.p_copy;        /* version-chain walk      */
%         } while (p_copy);
%     }
%     return p_act;
%
% So: >= 2 loads, and NOT a constant. Two facts matter and both were verified,
% not inferred from the RLU row (the comment above says why that inference is
% unsafe -- and it would have been WRONG here in the opposite direction):
%
%   1. The header read is UNCONDITIONAL. There is no counterpart to RLU's
%      `if (self->is_check_locks == 0)` pure-reader bypass. Multi-versioning
%      costs RLU's pure-reader fast path -- MV-RLU's reader always touches the
%      object. This is a real difference BETWEEN the two RLU rows.
%   2. The version-chain walk is length-dependent with no iteration limit, so
%      the cell is a floor, not a typical case. That is why it reads ">= 2"
%      rather than a number.
%
% BRANCH COLUMN -- from the same sources quoted above, all fast paths:
%   existence  if (likely(!egi)) return 1;                     -> 1, hinted
%   RLU        if (rlu_likely(self->is_check_locks == 0))      -> 1, hinted
%   MV-RLU     if (unlikely(!obj)) / is_obj_actual() inside
%              get_act_obj() / if (unlikely(p_copy))           -> >= 3, + loop
%   this work  urcu_txn_sw_list_resolve(), rcu-txn-sw-list.h:
%                 uintptr_t v = (uintptr_t) ptr;
%                 if (caa_unlikely(v & URCU_TXN_SW_LIST_PROXY_TAG)) {...}
%                 return ptr;                                  -> 1, hinted
%
% The point of carrying this column at all: EVERY design branches, ours
% included, so the branch is NOT where the comparison is decided and claiming
% otherwise would be unfair in our favour. The difference is that our test's
% operand is already in a register while everyone else must load theirs first.
% Being scrupulous here strengthens the argument rather than weakening it --
% do not drop the column to make the row look cleaner. Note also that our 1 is
% deliberately NOT bolded: bold marks where this design wins, and it does not
% win here.
%
% Deliberately NOT claimed: mvrlu_deref also does self->num_act_obj++ and
% self->num_deref++ (per-thread STORES on the read path). Those read like
% statistics counters and may be compiled out in a production build; we did not
% verify that, so we do not count them. Counting them would be exactly the kind
% of flattering-to-us overclaim the RLU note above exists to prevent.
%
% RLU HAS NO TLS. There is no __thread anywhere in rlu.h/rlu.c: `self` is a
% caller-held pointer parameter, so self->is_check_locks is a plain load through
% a register. Do not write "thread-local flag" -- it is a field of a per-thread
% STRUCT, which is not the same claim.
%
% Our own benchmark is what puts `self` in TLS (bench_list_scale.c:1505,
% `static __thread rlu_thread_data_t *tls_rlu`), fetched once per operation via
% rlu_self() at :1579 and then held in a register across the traversal. So it is
% hoisted out of the per-dereference path this table is about.
%
% DO NOT argue TLS access models (initial-exec vs global-dynamic under dlopen)
% against RLU. It is a property of how an embedder stores `self`, not of RLU's
% design -- RLU takes `self` as a parameter precisely to leave that open -- and
% it is amortised per operation, not per dereference. It also cuts both ways:
% rcu_read_lock() touches URCU_TLS(rcu_reader) under the same rules (urcu ships
% tls-compat.h for exactly this), so the point nets to ~zero and reads as
% reaching. The argument is strongest on ground where the difference is
% structural: existence and RLU must put their mechanism somewhere the reader
% can LOAD it; ours is already in the register.
% ---------------------------------------------------------------------------

\subsection{What each design charges the reader}
\label{sec:costwalk}

\paragraph{Every design tests something; only this one reaches the test for
free.} The difference is structural rather than a matter of constant factors,
but it is narrower than "some designs check the read path and this one does
not". Every row of \cref{tab:comparison} checks it, this one included. What
differs is what the check costs to reach its operand.

Existence's fast path is one acquire load of a per-element field, then a branch
on it. RLU's fast path, for a thread with no writes of its own, never touches
the object header --- but it still loads a flag from its per-thread state, once
per dereference, and branches on that. MV-RLU reads the object header
unconditionally and branches, then may walk a version chain. None of these loads
is expensive in isolation; all of them are memory references, and in each case
the branch cannot resolve until its load does.

A \pseudotxn's marker occupies tag bits of the slot, and those bits cost nothing
to occupy: a pointer to an aligned object leaves its low bits clear by
construction, so every transacted slot already wastes them, and the marker uses
only what alignment left idle (\cref{sec:alignment} sets out how many).
Recognizing the marker is then
\begin{center}
\code{((uintptr\_t) v \& tag) == tag}
\end{center}
where \code{v} is the value the reader's \code{rcu\_dereference} already
returned. That is a mask and a test on a register the reader was obliged to hold
anyway: no additional load and no additional cache line.

It is not free, and the residue is worth naming rather than rounding away. The
test is a conditional branch --- a couple of instructions, and one entry in the
branch predictor. It is hinted unlikely and is correctly predicted not-taken
whenever no commit is in flight, which is to say almost always; but a predicted
branch is not a free branch, and a reader in a tight traversal pays it on every
hop. So \cref{tab:comparison}'s \textbf{0} is zero \emph{loads}, not zero
instructions. And when the branch is taken---a commit is in flight on the slot
the reader just loaded---the reader does pay: it resolves the slot through the
proxy record and then the group's shared selector, a pair of loads the fast
path avoids and either of which may miss cache. That cost is bounded---only the
slots a commit touches, and only while it is in flight---but it is real, so
\cref{tab:comparison}'s figures describe the steady-state fast path.

That is the honest shape of the claim, and it is still the shape we want: the
branch is common to every design in the table, so it is not where the comparison
is decided. What separates this design is that its test operates on a value the
reader already had, while every other design must first go and fetch the thing
it will test --- from a per-element field, a per-thread flag, or an object
header. The marker rides the dereference the reader had to perform regardless.

\paragraph{One cost the table cannot express: comparing two pointers.} Counting
per dereference is the right unit for \cref{tab:comparison}, but it hides
something a traversal cannot avoid. A walk must also decide when it has
returned to where it started, and under RLU that comparison is not free: either
operand may be a copy rather than the object itself, so the two cannot be
compared directly and the walk must ask the library, which resolves both sides
to their originals --- reading an object header for each --- in an out-of-line
call, once per iteration. Measured on the list of \cref{sec:swlist}, that
comparison costs 5.5 loads and 6.0 branches per node visit, considerably more
than the dereference it accompanies. Substituting a plain pointer comparison,
which is sound only because no writer is running, removes exactly that: RLU's
dereference then costs the one extra load \cref{tab:comparison} states.

The table is therefore right about RLU's dereference and silent about RLU's
traversal, and the silence favors RLU over the \pseudotxn. The facility has no
such cost, for the same reason it has no per-element state: a resolved slot is
the node itself, so a reader compares two pointers with \code{!=}. The general point
is that a design which lets an object be represented by more than one address
must pay to establish identity, everywhere identity is needed --- and a
traversal needs it on every step.

\paragraph{The unconditional cost is space.} Instruction counts are the smaller
half of the argument, because the state is the part that cannot be made
conditional. RLU's header is prepended to every object and existence's field
sits in every element, for as long as the structure lives, whether or not any
mutation is in flight. That is a permanent addition to the footprint of the data
the readers traverse --- and footprint is a read-side cost, paid in cache rather
than in instructions. A \pseudotxn adds no per-element state at all, so in the
steady state both the structure a reader traverses and the path it takes through
it are byte-for-byte what they would be with no facility present.
And the cost is \emph{density} rather than size: what matters is how much of
each cache line the reader fetches is data it actually reads. Existence's own
structure makes the point, because the effect can be measured on one algorithm
by moving its fields. Its reference layout places the per-element header at
offset 16, which pushes the element's key out to offset 128 --- a second cache
line --- so a traversal touches two lines per element. Reordering so that the
link pointers, the key and the existence word share the first line, at
\emph{identical} size and semantics, is worth a factor of 1.9 to 2.2 in read
throughput. Segregating the state the reader never touches --- only the tagged
word is read during traversal; the group linkage, the reclamation head, the
callbacks and the writer's lock are not --- takes the element from 136 bytes to
32 and gains a factor of 2.1 to 2.5 over the reference layout.

That progression is why \cref{tab:readclass} reports existence at 1.14--1.21
rather than the 2.4--2.9 its reference layout yields: the larger number
measures a struct layout, not the design. What survives the tuning is the part
that cannot be tuned away --- one word per element, which is 8 bytes, which is
what takes a 24-byte element to 32 and puts fewer of them in each line. A
\pseudotxn has no such word to place well or badly.

% -----------------------------------------------------------------------
% NOTE (drafting; retained deliberately -- see README framing rule).
% The paragraph above REPLACES the older plan to cite the rcu_head
% segregated-vs-inline experiment (~2x at 1000 nodes) as a proxy for this
% effect. We no longer need a proxy: scripts/run_existence_layout.sh measures
% the SAME algorithm under three per-element layouts, so the spread is layout
% and nothing else (scripts/existence_layout.csv, engine b3e23f9f):
%
%   Mvisits/s, mean of 2      1 rdr        32       192
%   perfbook  136 B, 3 lines    326    10 311    50 205
%   packed    136 B, 1 line     623    20 023   108 942
%   split      32 B, 1 line     692    22 230   123 077
%   txn_sw_list 24 B            778    25 025   145 772
%
% The old caution still applies in spirit and must not be dropped: present
% this as evidence for the MECHANISM. It IS now a direct measurement of
% existence -- but of OUR builds of it, so the honest framing is what the
% concept costs when implemented well, which is why tab:readclass uses the
% split arm and says so in its caption. Do NOT quote the perfbook column as
% "existence's read cost".
% -----------------------------------------------------------------------

\paragraph{A pseudo-transaction charges the reader no memory.} The facility adds
no per-element state at all. Its marker lives in tag bits of the slot itself ---
the very word the reader had to load in order to proceed --- so recognizing it
costs no additional load and no additional line, only the predicted branch
accounted for above. And because the marker is transient, the structure a reader
traverses in the steady state is byte-for-byte the structure it would traverse
with no facility present.

\subsection{Transactional memory}
\label{sec:relstm}

Software transactional memory \citep{shavit1995stm, herlihy2003dstm} and its
object-based descendants \citep{fraser2004lockfreedom} offer the strictly
stronger contract this facility declines: a transaction is atomic \emph{and}
isolated, so a reader's whole read set is mutually consistent. Opacity
\citep{guerraoui2008opacity} names that reader-side guarantee precisely, and
naming it is what lets us be precise about what we do not provide.

\Cref{sec:fssense} argued that the divergence is not a matter of degree. STM
means the database sense of \emph{transaction}, which bundles isolation into the
word; we mean the journaling sense, which never has. The two are different
contracts, not stronger and weaker versions of one.

What we do take from STM is on the write side. Read-your-own-writes
(\cref{sec:swcomposition}) is its standard handling of a transaction that reads
a location it has itself already written, and we adopt it unchanged. It reaches
less far here: STM instruments every read, so a search made inside a transaction
is reconciled too, whereas ours reconciles only the slots a writer records and
leaves the traversal as plain \rcu. That shortfall is an obligation on the
embedder rather than a cost on the reader, which is the trade the whole facility
makes.

Nor is STM the only place that reader-side guarantee has been reached for. The
generation-number patents of \cref{sec:relrcu} give an \rcu reader a
version-consistent search: it stamps itself with a generation number and is then
answered as of that stamp, across every element it visits, however many updaters
have published since. That is not opacity --- there is no read set, no
validation, and no abort --- but it is a multi-element snapshot, and it is
therefore something this facility declines to offer and that line of work does
not. It is worth saying plainly, because the contrast in \cref{sec:fssense} is
between this facility and STM, and a reader could take from it that snapshots
are what one buys by leaving \rcu. They are not. They can be had within \rcu,
and the price is paid where \cref{sec:whyweakening} says it always is: on the
element and on the reader's path.

Transactional boosting \citep{herlihy2008boosting} is the closest relative on
the composition axis: coarse abstract locks over cheap physical operations,
rather than instrumented reads. It is the pattern a lock-plus-\fliplatch design
converges on, and the companion paper takes it up.

\subsection{Multi-pointer atomicity under RCU}
\label{sec:relrcu}

\rcu \citep{mckenney1998rcu, desnoyers2012urcu, mckenney2023perfbook} gives
readers a free ride and gives writers one atomic pointer store. Extending it to
several pointers at once is an old problem with three established answers.
\Cref{tab:comparison} places the two whose reader-side costs we could establish
from running code; the third is disclosed in patents and is treated here on its
claims.

\textbf{Existence structures} \citep{mckenney2015issaquah,
mckenney2016beyondissaquah, mckenney2026perfbookcode} give each element a field
naming an \emph{existence group}, and flip the group's state word to change the
membership of every element in it at once. The mechanism is the closest prior
art to this paper in what it asks of the reader: it is a shared word whose single
store switches a set's visibility to \rcu readers, followed by reclamation after
a grace period. The
technique has been public since at least January 2015 and was explained at
length in 2016.

\textbf{RLU} \citep{matveev2015rlu} gives each object a header able to name a
newer copy, logs a writer's copies, and makes them visible by advancing a global
clock. \textbf{MV-RLU} \citep{kim2019mvrlu} adds multi-versioning to relieve
RLU's synchronous writeback. Both provide a \emph{coherent snapshot} --- which
is strictly stronger than \cref{prop:monotone}, and must be said plainly: a
comparison of \emph{write} throughput between this facility and RLU would cross
guarantee classes and mean little. What \cref{tab:readclass} does compare is the
read side, where the crossing is informative precisely because it runs against
us: RLU charges the reader more and delivers more, and the gap is the price of
the snapshot. Against existence structures the classes are nearer still and the
comparison is cleaner.

\textbf{Generation-numbered versions} are the third answer and the oldest,
disclosed in a family of patents naming McKenney and assigned to IBM
\citep{mckenney2008cyclicsearch511, mckenney2010cyclicsearch082,
mckenney2011cyclicparallel778, mckenney2014rcucyclic535,
mckenney2015rcucyclic803}. An updater builds new versions of the elements it is
changing, stamps them with a generation number, links each to its prior version
in both directions, splices them in so readers can reach them, and then advances
a single global generation number; prior versions and their links are freed
after a grace period. A reader stamps its search with the global generation
number it starts from and, on meeting an element stamped differently, follows
the version links to the version matching its own stamp. The parallel-updater
member \citep{mckenney2011cyclicparallel778} hands generation numbers out in the
order updaters begin and advances the global number only across the completed
prefix, so no reader observes a partially published later update; the later two
\citep{mckenney2014rcucyclic535, mckenney2015rcucyclic803} make the
pointer-and-generation pair a single word one store can swing.

That last step is a shared word whose single store makes a set of new versions
current, which is the same shape as the flip of \cref{sec:fliplatch}, and it
predates the existence-structure disclosures by roughly a decade
(\cref{sec:relclaims}). The two designs diverge on what the reader does with it.
An existence flip and a \fliplatch both switch the reader between two states and
leave the old one unreachable; the generation-number scheme keeps old and new
versions simultaneously reachable and asks the reader to select among them by
comparing stamps and walking version links. That is what buys its readers a
version-consistent search --- and what it costs is exactly the axis
\cref{sec:whyweakening} argues over: a generation number on the element, version
links to traverse, and a reader that must compare and may walk rather than
resolve and proceed.

% ===========================================================================
\section{Related Work}
\label{sec:related}

\subsection{Publishing is not compare-and-swapping}
\label{sec:relmcas}

Multi-word compare-and-swap is the obvious neighbour and the wrong comparison,
so it is worth separating carefully.

Harris, Fraser and Pratt's practical \mcas \citep{harris2002mcas} makes $k$
words change together under concurrent writers: a descriptor names the records,
threads that meet it \emph{help} drive it, and RDCSS resolves the ordering
hazard between installing a descriptor and reading it. Fraser's thesis and the
journal generalization develop it into a lock-free toolkit
\citep{fraser2004lockfreedom, fraser2007locks}; wait-free constructions follow
\citep{sundell2011waitfreemcas, feldman2015waitfreemcas}; descriptor reuse
\citep{arbelraviv2017descriptors} and lower-overhead \kcas
\citep{guerraoui2020mcas} refine the cost.

The \fliplatch is not one of these. It performs \emph{no compare-and-swap at
all}: under exclusion there is nothing to compare against, no contender to lose
a race to, and no helping to arrange. It solves the narrower problem of atomic
multi-slot \emph{publication} --- making a prepared set of values visible at
once to readers who are not writers --- rather than atomic multi-slot
\emph{update} among competing writers.

What it does inherit from that line is the shape: a descriptor of records plus a
status word whose transition decides the outcome, and resolution routed through
the status rather than through the slot's plain value. That shape is Harris's,
and we claim none of it. The companion paper, which does contend, engages this
work properly.


\subsection{Patented prior art}
\label{sec:relpatents}

The atomic-move-under-\rcu problem has been the subject of a series of patents
assigned to IBM --- most naming McKenney, one Triplett --- and they are prior art
whether or not any remains in force. The generation-number family is treated as a
technical answer in \cref{sec:relrcu}; four others bound what this paper may
claim, and are taken in turn here.

\citet{mckenney2008filerename926} (filed 2004) renames a file while permitting
lock-free lookups by linking a \emph{temporary record} into the hash chain and
removing it afterwards --- a transient object, erected for the mutation and gone
after. Its readers, however, \emph{wait}: the claim recites that the temporary
record ``causes the look-ups to wait'' until the rename completes.
\citet{triplett2010keyupdate851} (filed 2006) moves an element between lists by
making it ``appear simultaneously'' in both under one key and then changing the
key --- appear-in-both, belong-to-one, which is the existence notion in embryo.
\citet{mckenney2018atomicmove907} (filed 2015) atomically updates a
status-indicating entity to transfer an element's validity from an original to a
copy. \citet{mckenney2019atomicmove638} (filed 2016) changes an existence
group's state field ``using a single store operation'' so that, where there are
plural elements, one set is deemed to exist and another deemed not to.

% -----------------------------------------------------------------------
% NOTE (drafting; retained deliberately -- see README framing rule).
% Cite the patents as PRIOR ART and stop there. Do not editorialise about
% legal status. All ten in the bib are recorded as lapsed or expired as of
% 2026-07-16, and six of those are fee-lapsed with term still running, which
% under 37 CFR 1.378 can be reinstated -- so "lapsed" is a point-in-time
% observation, not a fact about the future, and it is not this paper's
% business either way. The citation's job is to date the disclosure, not to
% assess enforceability. Saying less here is both safer and more correct.
% -----------------------------------------------------------------------

\subsection{What this paper claims}
\label{sec:relclaims}

Given the above, the claim has to be stated as a conjunction, because each
ingredient separately is old.

The \emph{flip} is not new: a single store switching a set's visibility to \rcu
readers is \citet{mckenney2019atomicmove638} and was public at CppCon in 2016
\citep{mckenney2016beyondissaquah}. It is older than that. A single store to a
shared word that makes a set of independently prepared new versions current is
the generation-number scheme of \citet{mckenney2008cyclicsearch511}, filed in
2004, and the parallel-updater member \citep{mckenney2011cyclicparallel778} is
already a rule for gating that store so that a set becomes current in one step.
What those readers do after the store differs (\cref{sec:relrcu}); the store
itself does not. The \emph{transient object} is not new: a
per-mutation record, removed afterwards, is \citet{mckenney2008filerename926}
from 2004. \emph{Pointer tagging} is not new by any measure. Hardware tag bits on
pointer-sized words shipped in a commercial platform with the IBM System/38
\citep{levy1984capability}, and the software practice this facility actually
uses --- stealing the low bits that alignment guarantees --- was old enough by
1993 to be surveyed as folklore that implementers had been absorbing piecemeal
or re-inventing \citep{gudeman1993typeinfo}. Nearer to hand, existence
structures tag bit~0 of their own membership field. The \emph{descriptor-and-status shape} is Harris's
\citep{harris2002mcas}.

What we have not found in that prior art is the combination: a marker that is
\textbf{transient} (so the structure's steady state is unchanged),
\textbf{non-blocking for readers} (they resolve through it to old-or-new rather
than waiting on it, unlike \citet{mckenney2008filerename926}), and
\textbf{free of per-element state} (it lives in the slot's spare bits rather
than in a field or header, unlike existence and RLU) --- together with the
consequence that follows: because the marker is per-slot rather than per-object,
a \pseudotxn can span slots whose encodings differ, which forces the per-record
tag (\cref{sec:pertag}) and which we are likewise not aware of in prior work.
% -----------------------------------------------------------------------
% NOTE (drafting; retained deliberately -- see README framing rule).
% "We have not found" and "we are not aware of" are the right register --
% they are honest, and they are what a careful reader will believe. "Novel"
% and "the first" are neither, and the prior-art sweep behind this section
% (10 patents, plus the talks) is exactly what makes the weaker phrasing
% credible rather than evasive. Do not let an editing pass strengthen these
% into claims.
% -----------------------------------------------------------------------

% ===========================================================================
\section{Conclusion}
\label{sec:conclusion}

\rcu gives a writer one atomic operation, and the limit is not academic. It is visible
in the interface of a shipped container: Userspace \rcu's list maintains a
backward chain that no reader is allowed to walk, and the missing reverse
iterator is not an oversight but the thing that licenses the plain store to
\code{prev} (\cref{sec:onestore}). Where an update needs two edges to move
together, the remedies have been to copy a larger region, to wait out a grace
period, to impose memory ordering at a read-side cost, or to withdraw the
operation.

We have described a fifth answer that leaves the contract intact. A \pseudotxn
is a frozen set of \record{slot}{old}{new} records made visible by one commit
step; it does not add a phase to \rcu's build--publish--reclaim discipline but
widens the middle one, so that a single store moves a set of slots rather than a
single slot (\cref{sec:buildpublish}). The \fliplatch realizes it for a single
writer with no compare-and-swap, no conflict detection and no abort: every proxy
in a group names one shared selector word, and one release store switches the
whole set (\cref{sec:fliplatch}). The marker that makes this possible rides in
spare bits of the slot the reader had to load anyway, so it costs no per-element
state and no additional load, and it is gone once the mutation settles
(\cref{sec:tags}).

The guarantee is deliberately weaker than transactional memory, and we have
tried to be more precise about the weakness than about the strength. The write
set is atomic; a reader is not a transaction of any kind, and a traversal
spanning several
slots may straddle a commit. There is no opacity and no multi-read snapshot.
What a reader gets instead is monotonicity --- never new-then-old
(\cref{sec:monotone}) --- and that middle ground is exploitable rather than
merely tolerable: an embedder that commits in the reverse of its readers'
traversal order gets old-prefix, new-suffix observations and needs no
per-traversal snapshot at all (\cref{sec:publishorder}). The weakening is
targeted.\footnote{The sentiment is not ours, and neither is the better title
for it: \citet{alglave2013weakness}, of weak memory models rather than of
transactional guarantees.} Short of copying the structure outright --- which is the cost that
made the operation expensive to begin with --- isolation is a charge on the
reader, and the reader is what \rcu exists to keep free.

The costs are real and we have stated them where they fall rather than where
they flatter. Exclusion is a precondition of the mechanism, not a caveat on it,
and violating it corrupts silently (\cref{sec:exclusion}). Naming individual
slots means a commit's width grows with the number of edges an operation
re-points, which is the axis on which a design that collapses membership into a
single object will win (\cref{sec:costcomparison}). The tag contract constrains
every value a transacted slot may hold, and an embedder that breaks it does not
get a torn read but a wild pointer (\cref{sec:tagcontract}).

What we claim is a conjunction rather than a component. The marker is transient,
non-blocking for readers, and free of per-element state; each of those exists
separately in prior work, and we have not found them together, nor found the
per-record tag that having them together forces (\cref{sec:related}). We have
described the mechanism, and the variants we did not take
(\cref{sec:variations}), in enough detail to be reimplemented from this paper
alone --- a description that cannot be built from is not a disclosure.

\paragraph{The rest of the series.} This paper is the \sw case, and it is
deliberately the first: the \fliplatch is the model at its smallest, where a
commit is one store and the hard parts are contracts rather than machinery. A
companion paper takes up the \mw facility, which realizes the same model for
concurrent writers, paying in machinery this paper has no need of for the
exclusion it drops. Because the two share the model of \cref{sec:model} and the
status vocabulary, an embedder that outgrows exclusion changes which facility it
links, not how it reasons. The programming model the two facilities share, and the
evaluation, follow in their own papers.

% ===========================================================================
\section*{Availability}

% ---------------------------------------------------------------------------
% Pin the COMMIT, never the branch. urcu-txn-dev is a development tip and will
% move; the disclosure is the tree as of this specific commit and nothing else.
% For a defensive publication the dated, specific version IS the deliverable,
% so a reader must be able to fetch exactly what this paper describes years
% from now, without depending on where development happened to go next.
% ---------------------------------------------------------------------------
The facility described here is free software, distributed under
\code{LGPL-2.1-or-later} as part of Userspace \rcu. Every mechanism this paper
describes is present in commit
\code{b3e23f9f558ce895aebe300ddb33aee9281137dc} (2026-07-25), on branch
\code{urcu-txn-dev} at \url{https://github.com/compudj/userspace-rcu-dev}. The
commit hash, not the branch, identifies the disclosed version: the branch is a
development tip and will move.

Within that tree, the material of this paper is the following headers, all under
\code{include/urcu/}:

\begin{description}\setlength{\itemsep}{2pt}
\item[\code{rcu-txn-sw.h}] The \fliplatch: flip groups, proxies, resolution, and
      the four-step lifecycle of \cref{sec:lifecycle}. Its ordering note is the
      source of \cref{sec:monotone} and \cref{sec:publishorder}, and its
      read-your-own-writes pair (\code{urcu\_txn\_sw\_load} /
      \code{urcu\_txn\_sw\_record\_chain}) is the composition mechanism of
      \cref{sec:swcomposition}.
\item[\code{rcu-txn-sw-list.h}] The bidirectional list of \cref{sec:wrappers},
      whose two chains stay mutually coherent.
\item[\code{rcu-txn-sw-hlist.h}] The hash list.
\item[\code{rcu-txn-status.h}] The commit-status enum
      (\cref{sec:model}) that both front-ends return, so an embedder tests a
      commit the same way whichever facility it drives.
\end{description}

\noindent Three further headers are named here for orientation rather than as
contributions of this paper: \code{rcu-txn-mcas.h} is the shared multi-slot
commit primitive --- one descriptor whose slots install either by
compare-and-swap, the \mw discipline the companion paper takes up, or by the
exclusive park this paper uses --- \code{rcu-txn.h} is the transaction front-end
over it, and
\code{rcu-txn-sw-bitmap.h} applies the same tag contract to a non-pointer word
--- outside this paper's pointer-focused scope.

This paper is licensed CC BY 4.0. It is published to place the design, and the
alternatives of \cref{sec:variations}, in the public record as prior art, so
that they remain free for anyone to implement and use.

\section*{Acknowledgements}

The author thanks Paul E. McKenney for his careful review of this article.

\bibliographystyle{plainnat}
\bibliography{../common/urcu-txn}

\end{document}
